\section{Introduction}\label{sec:Introduction}

% =============================================================================
% INTRODUCTION STRUCTURE (Following OR/OM Conventions)
% =============================================================================
% Part 1: Opening paragraphs - Problem context and motivation (~1 page)
% Part 2: Model setting and key insight (~0.5 page)
% Part 3: Approach and main results (~0.5 page)
% Part 4: Contributions (bulleted)
% Part 5: Roadmap
% Part 6: Literature review (subsection at end)
% =============================================================================

% =============================================================================
% PART 1: OPENING PARAGRAPHS - PROBLEM CONTEXT AND MOTIVATION
% =============================================================================

The rapid expansion of on-demand ride-hailing platforms has transformed urban transportation.
By algorithmically matching passengers with nearby drivers, platforms such as Uber, Lyft, and DiDi have reduced the search costs that traditionally plagued taxi markets.
As these platforms mature, they increasingly serve heterogeneous demand classes from a shared pool of drivers.
Such heterogeneity arises naturally: dynamic pricing creates segments with different willingness to pay, service tiers (e.g., UberX vs.\ Uber Comfort) attract distinct customer populations, and subscription programs differentiate members from casual users.
This heterogeneity introduces a tension.
Platform managers must balance revenue from high-margin customers against maintaining service quality for price-sensitive segments.
The shared driver pool means that prioritizing one class necessarily affects the other.

The operational complexity stems from a distinctive feature of ride-hailing: pickup times are endogenous.
Unlike traditional queues where service times are exogenous, the time required to pick up a matched passenger depends jointly on the spatial density of idle drivers and waiting passengers.
When demand surges or driver supply contracts, pickup times lengthen for all classes, leading to two forms of customer loss: passengers may abandon before being matched (pre-match abandonment), or cancel after matching if the quoted pickup time is too long (post-match cancellation).
The key challenge is that matching decisions interact across classes.
When the platform matches a passenger from one class, an idle driver is consumed, reducing the pool available for other classes.
This capacity spillover degrades pickup efficiency for the remaining classes through the same spatial mechanism.
Consequently, the optimal matching policy for one class depends on the state and policy of the other.
Static priority rules or decoupled optimization methods do not capture this interdependence.

This paper develops an analytical framework for coordinated matching control.
We use pickup rate thresholds as control levers: for each class, the platform specifies a minimum pickup rate below which matches are deferred.
The central question is how to set these thresholds optimally, accounting for the coupling between classes.

% =============================================================================
% PART 2: MODEL SETTING AND KEY INSIGHT
% =============================================================================

We consider a ride-hailing system with two demand classes sharing a pool of drivers.
Passengers arrive according to class-specific Poisson processes and join separate queues.
The matching rate for each class follows a Cobb-Douglas spatial production function, where the rate increases in both the number of waiting passengers and the count of idle drivers.
This formulation, established in prior work on single-class systems, captures the economies of spatial density: more passengers and more drivers lead to faster matching.
Passengers exhibit two-stage impatience: they may abandon while waiting in queue, or cancel after matching if the pickup takes too long.
The platform controls matching through pickup rate thresholds: a passenger is matched only when the current pickup rate exceeds a class-specific threshold.

Our main theoretical contribution is the derivation of the \emph{Key Matching Indices} (KMI), denoted $(\zeta_1, \zeta_2)$, which provide observable criteria for optimal threshold selection.
Each index decomposes into three interpretable terms:
(i) \emph{abandonment reduction}, the benefit from reduced queue length and shorter waits;
(ii) \emph{direct opportunity cost}, the cost of consuming an idle driver; and
(iii) a \emph{cross-market term}, the effect on the other class through the shared driver pool.
The third term is crucial: it captures the capacity spillover that single-class models and decoupled methods ignore.
We prove that optimal thresholds are characterized by driving both indices to unity: $(\zeta_1, \zeta_2) = (1,1)$.
This condition holds for both welfare maximization and revenue maximization, though the precise form of the indices differs between objectives.

% =============================================================================
% PART 3: APPROACH AND MAIN RESULTS
% =============================================================================

Building on the single-class spatial framework of \citet{WangZhangHu2024}, we extend the model to heterogeneous demand and derive fluid approximations that remain tractable for optimization.
The steady-state analysis yields the KMI characterization, which reduces the high-dimensional optimization problem to a system of coupled nonlinear equations.
Based on the KMI structure, we design a coordinated control algorithm that uses the Jacobian of the KMI map to adjust thresholds dynamically.
Because the cross-market terms create genuine coupling, the Jacobian captures capacity spillovers that decoupled gradient methods miss.
Numerical experiments demonstrate that the algorithm converges to optimal thresholds from any initial state.
We systematically characterize the welfare-revenue frontier across different load conditions.
Under low load, both objectives can be optimized simultaneously with generous thresholds.
Under congestion, however, revenue maximization requires disproportionate rationing of the regular class to preserve service quality for premium customers.
The KMI framework quantifies this trade-off and explains why heuristic tuning methods perform poorly in congested conditions.

% =============================================================================
% PART 4: CONTRIBUTIONS
% =============================================================================

\subsection{Contributions}

Our main contributions are as follows:

\begin{enumerate}
    \item We derive Key Matching Indices (KMI), $(\zeta_1, \zeta_2)$, which are observable, state-dependent indicators for optimal threshold selection.
    Each index decomposes into three interpretable terms: abandonment reduction, direct opportunity cost, and a cross-market term that captures how each class's threshold affects the other.
    We prove that the condition $(\zeta_1, \zeta_2) = (1,1)$ characterizes optimal thresholds for both welfare and revenue objectives, reducing steady-state optimization to a tractable system of coupled equations.
    The welfare and revenue indices are structurally distinct: the former measures opportunity costs in physical units (throughput per hour), the latter in value units (revenue per hour), yet the unity condition holds for both.

    \item We design a coordinated control algorithm that uses the Jacobian of the KMI map to adjust thresholds dynamically.
    The Jacobian is diagonally dominant (each threshold's own effect exceeds its cross-market effect), which ensures invertibility and local stability of the Newton-type iteration.
    By internalizing the cross-market coupling, the algorithm converges to optimal thresholds from any initial state and outperforms decoupled or myopic tuning rules that ignore cross-class externalities.

    \item We characterize the welfare-revenue frontier across load conditions.
    While both objectives align under low load, they diverge significantly under congestion, where revenue maximization requires disproportionate rationing of the regular class to preserve spatial density for premium users.
    Specifically, the revenue-optimal policy reduces the premium queue by over $40\%$ relative to the welfare-optimal policy but more than doubles the regular queue.
    The KMI framework reveals asymmetric sensitivities that explain why heuristic tuning methods perform poorly in congested regimes.

    \item Relative to static priority rules and decoupled heuristics, our KMI-based control delivers substantial performance gains, particularly under congestion: the welfare-optimal KMI policy improves welfare by $19\%$ and revenue by $9\%$ over static dispatch.
    The primary source of these gains is the internalization of cross-market externalities that decoupled policies ignore.
    From a managerial perspective, the results support a shift from static or heuristic dispatch rules toward model-based threshold management, with the KMI serving as real-time indicators for service differentiation.
\end{enumerate}

% =============================================================================
% PART 5: LITERATURE REVIEW
% =============================================================================

\subsection{Related Literature}

Our work connects to four streams of research: spatial matching models, service systems with abandonment, dynamic matching and priority mechanisms, and platform objectives.

\paragraph{Spatial Matching and Endogenous Pickup Times.}
The defining feature of ride-hailing is that pickup times are endogenous, depending on the spatial density of drivers and passengers.
This concept has been rigorously formalized by \citet{BesbesCastroLobel2021, BesbesCastroLobel2022}, who establish key results for single-class settings.
Our work builds directly on the stochastic-fluid framework of \citet{WangZhangHu2024}, which characterizes matching rates via Cobb-Douglas production functions and derives optimal thresholds under abandonment.
Building on this foundation, we extend the framework to heterogeneous demand, where capacity spillovers arise from the shared driver pool that is absent in single-class models.
An influential stream of recent research incorporates strategic driver behavior into spatial matching models.
The seminal work of \citet{doi:10.1287/msom.2023.1221} develops a rigorous game-theoretic fluid model that captures the interplay between platform control and self-interested driver decisions, yielding a key insight: platforms may optimally reject demand at low-demand locations to induce strategic repositioning toward high-demand areas.
Their characterization of equilibrium behavior under different control regimes provides essential benchmarks for evaluating platform policies.
\citet{GargNazerzadeh2022} further explore driver strategic responses, while other extensions include pricing for matching failures \citep{Castillo2025, CastilloKnoepfleWeyl2017} and routing considerations \citep{BravermanDaiLiuYing2019}.
Complementary to this line of work, which focuses on spatial repositioning across locations, our work examines how threshold-based matching controls can manage service differentiation across heterogeneous customer classes sharing a common driver pool.

\paragraph{Service Systems with Abandonment.}
Customer impatience is a central feature of on-demand services.
Fluid approximations for many-server queues with abandonment began to emerge following the work of \citet{whitt_FluModMul_2006}; see also \citet{Zhang2013, LiuWhitt2012, LiuWhitt2014, Kang2014, LongZhang2014} and the survey by \citet{dai_ManSerQue_2012}.
\citet{bassamboo_AccFluMod_2010} further show that fluid-based capacity prescriptions achieve $O(1)$ accuracy in the presence of customer abandonment.
\citet{doi:10.1287/mnsc.2023.01513} derive scaling laws for the cost of impatience in dynamic matching, while \citet{Bai2019} and \citet{feng2021we} incorporate waiting time into customer utility.
Building on these models, our work provides stochastic integral equations that capture two-stage customer loss (pre-match abandonment and post-match cancellation) within a spatial priority framework, yielding tractable fluid limits for optimization.

\paragraph{Dynamic Matching and Priority Mechanisms.}
Dynamic matching networks constitute an active research area.
\citet{GurvichWard2014} prove the asymptotic optimality of discrete review matching policies under large market assumptions, while \citet{OzkanWard2020} develop revenue-maximizing dispatch control through fluid limit analysis.
In a related many-server setting, \citet{bassamboo_DynRouAdm_2005} establish the asymptotic optimality of joint admission and routing control using fluid limits.
\citet{HuZhou2022} analyze market thickness in dynamic type matching, and further studies explore market thickness effects \citep{doi:10.1086/704761, baccara2020optimal}, shared-ride efficiency \citep{Taylor2024}, and two-stage stochastic matching \citep{CandoganBalseiro2024}.
Index-based priority rules have a rich history in multi-class queueing.
The generalized $c\mu$ rule \citep{vanmieghemDynamicSchedulingConvex1995, mandelbaumSchedulingFlexibleServers2004} assigns priority through a product of service rate and marginal delay cost; \citet{longDynamicSchedulingMulticlass2020} extend this to many-server queues with abandonment via a $c\mu/h$ rule that divides by the patience hazard rate, building on the virtual allocation framework of \citet{long_VirAllPol_2019}.
Related work on priority scheduling in overloaded systems includes \citet{PuhaWard2019, AKS2013, MandelbaumPats1995, MandelbaumPats1998}.
Our approach implements priority through state-dependent pickup rate thresholds rather than discrete class prioritization.
Like the $c\mu/h$ index, each KMI incorporates abandonment costs, but it additionally includes a cross-market externality term that captures capacity spillovers through the shared driver pool.
\citet{doi:10.1287/msom.2022.0216} also employ prioritization mechanisms, though their focus is on peak/off-peak balancing whereas we address service differentiation across heterogeneous classes.

\paragraph{Platform Objectives and Economics.}
A substantial body of research examines how ride-hailing platforms balance revenue, welfare, and labor objectives.
Pricing and wage mechanisms have been extensively studied \citep{Taylor2018, Benjaafar2022, CHEN2022108329, hu2020price}, as have dynamic pricing with demand learning \citep{DenBoer2014, KeskinZeevi2014, BesbesZeevi2015} and surge pricing strategies \citep{CuiOrhunDuenyas2022, yan2020dynamic}.
\citet{OZKAN20201149} and \citet{BanerjeeRiquelmeJohari2015} analyze joint pricing and matching decisions, while \citet{BanerjeeFreundLykouris2022, IglesiasRossiZhangPavone2019, GurvichLariviereMoreno2019} study spatial capacity planning.
Complementary to this pricing-centric literature, we treat prices as exogenous and focus on matching control through threshold policies.
We employ the concept of marginal value conservation of idle drivers, inspired by state-dependent control approaches in closed queueing networks \citep{BanerjeeKanoriaQian2018, PowellSchultz2004}, to derive unified optimality criteria that apply to both welfare and revenue objectives.

% =============================================================================
% PART 6: ROADMAP
% =============================================================================

The rest of the paper is organized as follows.
Section~\ref{sec:Stochastic_Model} formulates the stochastic model for the two-class system with endogenous pickup times and two-stage customer loss.
We construct and analyze the corresponding fluid model in Section~\ref{sec:Fluid_Model}.
In Section~\ref{sec:matching_policies}, we derive the Key Matching Indices and characterize optimal thresholds for both welfare and revenue objectives.
Section~\ref{sec:numerical_experiments} presents numerical experiments demonstrating the performance of our coordinated control algorithm, and Section~\ref{sec:Conclusion} concludes with managerial insights and future research directions.
%%%%%%%%%%%%%%%%%%%%%%%%%%%%%%%%%%%%%%%%%%%%%%%%%%%%%%%%%%%%%%%%%%%%%%

\section{Stochastic Model}\label{sec:Stochastic_Model}

We formulate a stochastic model for a ride-hailing system serving two heterogeneous customer classes from a shared pool of drivers.
The key feature distinguishing this system from traditional queueing models is that pickup times are endogenous, depending on the spatial density of both waiting customers and idle drivers.
Waiting customers may abandon before being matched (pre-match abandonment), and even after matching, lengthy pickups can lead to cancellation (post-match cancellation).
The platform controls matching through threshold policies that specify minimum pickup rates required for each class.

The system operates with $K$ homogeneous drivers who transition among five states.
A driver can be idle and available for dispatch (count $Z_0(t)$ at time $t$), assigned to pick up a Class $i \in \{1,2\}$ customer (count $Z_{1i}(t)$), or actively transporting a Class $i$ customer (count $Z_{2i}(t)$), where the driver population is conserved: $Z_0(t) + \sum_{i=1}^2 Z_{1i}(t) + \sum_{i=1}^2 Z_{2i}(t) = K$ for all $t \geq 0$.

Customers form class-specific queues, with $Q_i(t)$ denoting the number of Class $i$ customers awaiting a match.
The two classes are differentiated by service tier (e.g., UberX versus Uber Comfort) or dynamic pricing, which creates customer segments with different willingness to pay.
What matters for our analysis is that both classes draw from the same idle-driver pool $Z_0(t)$, creating competition for scarce capacity.
This coupling through the shared driver pool is the key source of interaction between the two classes and will play a central role in our optimization analysis.

Class $i$ customers arrive according to a Poisson process $A_i(t)$ with rate $\lambda_i > 0$.
While waiting in queue, customers are impatient.
Each Class $i$ customer has an exponentially distributed patience with rate $\theta_{0i} \geq 0$.
If a customer's patience expires before being matched, they abandon.
The cumulative abandonment count is $R_{0i}(t) = W_{0i}(\int_0^t \theta_{0i} Q_i(s) ds)$, where $W_{0i}$ is an independent unit-rate Poisson process.

The defining feature of ride-hailing systems is that pickup times depend on spatial density.
When more customers and more drivers are present, the average distance between a matched pair is shorter, leading to faster pickups.
For a customer-driver pair matched at time $t$, we model the pickup time as a random variable following an exponential distribution with rate $\mu_{1i}(t)$.
Based on the spatial model introduced in \citet{WangZhangHu2024}, we make Assumption~\ref{assumption:spatial_model_rate}.

\begin{assumption}[Spatial Matching Function]
\label{assumption:spatial_model_rate}
When a Class $i$ customer is matched at time $t$, the pickup time for that match is exponentially distributed with rate
\begin{equation}
\mu_{1i}(t) = C (Q_i(t))^{\alpha_1} (Z_0(t))^{\alpha_2}, \quad i=1, 2, \label{eq:spatial_model_rate}
\end{equation}
where $C, \alpha_1, \alpha_2 > 0$ are constants.
Conditional on the system state at the moment of matching, the pickup time is independent of all other random variables.
\end{assumption}

We refer to $\mu_{1i}(t)$ as the pickup rate for Class $i$.
The parameters $C, \alpha_1, \alpha_2$ capture the spatial characteristics of the operating region.
Higher values of $C$ correspond to denser urban areas with shorter distances.
The exponents $\alpha_1$ and $\alpha_2$ reflect returns to density: more waiting customers increase the likelihood that one is located near an idle driver, and more idle drivers increase coverage.
A critical feature is that the rate $\mu_{1i}(t)$ is determined at the instant of matching and remains fixed throughout the subsequent pickup phase.

\paragraph{\bf Threshold-Based Matching Policy.}
To maintain service quality, the platform applies a threshold-based matching policy by imposing class-specific thresholds $\mu^*_{11}$ and $\mu^*_{12}$.
A Class $i$ customer is matched only when the pickup rate exceeds the threshold $\mu^*_{1i}$, i.e., when $\mu_{1i}(t) \geq \mu^*_{1i}$.
When $\mu_{1i}(t) < \mu^*_{1i}$, the platform defers matching and allows the number of waiting customers and idle drivers to accumulate until the pickup rate improves.
Define the eligibility indicator $E_i(t) = \mathbf{1}\{\mu_{1i}(t) \geq \mu^*_{1i}\}$ and the activity indicator $I_i(t) = E_i(t)\mathbf{1}\{Q_i(t)>0\}$.
Class $i$ is active at time $t$ if both eligibility and queue nonemptiness conditions hold.

\paragraph{\bf The Matching Process.}
Matches occur at event times triggered by either (i) customer arrivals or (ii) drivers becoming idle.
Since a matching can only happen when a customer arrives or a driver becomes idle, we denote $N_D(t) = \sum_{j=1}^2 R_{1j}(t) + \sum_{j=1}^2 D_{2j}(t)$ as the cumulative count of idle driver releases, where $R_{1j}$ and $D_{2j}$ are pickup cancellations and trip completions defined below.
We now describe the matching logic at these event times which leads $\mu_{1i}(t)$ to become bigger than the threshold $\mu^*_{1i}$.

When a Class $i$ customer arrives, a match is formed immediately if an idle driver is available ($Z_0(t^-) > 0$) and Class $i$ is eligible ($E_i(t^-) = 1$).
Otherwise, the customer joins the queue.

When a driver becomes idle, the matching logic depends on which classes are active at that moment:
\begin{itemize}
    \item If exactly one class $k$ is active ($I_k(t^-) = 1$ and $I_j(t^-) = 0$ for $j \neq k$), the driver is matched to Class $k$.
    \item If both classes are active ($I_1(t^-) = I_2(t^-) = 1$), the driver is assigned to Class 1 with probability $p_1(t^-) = Q_1(t^-)/(Q_1(t^-) + Q_2(t^-))$ and to Class 2 with probability $p_2(t^-) = Q_2(t^-)/(Q_1(t^-) + Q_2(t^-))$.
    \item If no class is active, the driver remains idle.
\end{itemize}

This matching rule leads to the following stochastic integral representations for cumulative matches:
\begin{align}
M_1(t) &= \int_0^t \mathbf{1}\{E_1(s^-), Z_0(s^-)>0\} dA_1(s) \notag \\
&\quad + \int_0^t \left[ I_1(s^-)(1-I_2(s^-)) + I_1(s^-)I_2(s^-) p_1(s^-) \right] dN_D(s), \label{eq:match_class1} \\
M_2(t) &= \int_0^t \mathbf{1}\{E_2(s^-), Z_0(s^-)>0\} dA_2(s) \notag \\
&\quad + \int_0^t \left[ I_2(s^-)(1-I_1(s^-)) + I_1(s^-)I_2(s^-) p_2(s^-) \right] dN_D(s). \label{eq:match_class2}
\end{align}
Here $s^-$ denotes the left limit, capturing the system state immediately before each event.

Once a Class $i$ customer is matched at time $\tau$, the assigned driver attempts pickup.
Two events compete: pickup completion (exponentially distributed with the rate $\mu_{1i}(\tau)$ fixed at matching) and customer cancellation (exponentially distributed with rate $\theta_{1i} \geq 0$).
If pickup completes first, the pair transitions to the trip phase (state $Z_{2i}$).
If the customer cancels first, the driver returns to the idle pool and the customer exits the system.

We denote cumulative pickup completions by $D_{1i}(t)$ and cumulative cancellations by $R_{1i}(t)$.
The construction of these processes accounts for pairs initially in the pickup state and pairs matched during $(0,t]$.
For the $k$-th matched pair with matching time $\tau_k$, let $Y_k \sim \text{Exp}(\mu_{1i}(\tau_k))$ be the pickup time and $V_k \sim \text{Exp}(\theta_{1i})$ be the remaining patience.
The cumulative outcomes are formally given by sums of indicator functions over all relevant pairs.

Trip durations are exponentially distributed with a class-independent rate $\mu_2 > 0$.
Upon trip completion $D_{2i}$, the driver returns to the idle state.
The cumulative trip completions evolve as $D_{2i}(t) = W_{2i}(\int_0^t \mu_2 Z_{2i}(s) ds)$, where $W_{2i}$ are independent unit-rate Poisson processes.

The queue dynamics follow from arrivals, abandonments, and matches:
\begin{equation}
Q_i(t) = Q_i(0) + A_i(t) - R_{0i}(t) - M_i(t), \quad i=1,2. \label{eq:Qi}
\end{equation}
Driver state transitions are governed by matching, pickup outcomes, and trip completions.
The idle pool evolves as
\begin{equation}
Z_0(t) = Z_0(0) + \sum_{j=1}^2 R_{1j}(t) + \sum_{j=1}^2 D_{2j}(t) - \sum_{j=1}^2 M_j(t). \label{eq:Z0}
\end{equation}
Pickup-state populations satisfy
\begin{equation}
Z_{1i}(t) = Z_{1i}(0) + M_i(t) - D_{1i}(t) - R_{1i}(t), \quad i=1,2, \label{eq:Z1i}
\end{equation}
and trip-state populations evolve as
\begin{equation}
Z_{2i}(t) = Z_{2i}(0) + D_{1i}(t) - D_{2i}(t), \quad i=1,2. \label{eq:Z2i}
\end{equation}
Together with the definitions of the cumulative processes, these flow balance equations fully specify the stochastic dynamics of the system.

In summary, we have introduced a stochastic model for a ride-hailing system with heterogeneous customer classes and two-stage customer loss (pre-match abandonment and post-match cancellation).
The primitive parameters of the system are $\phi = (K, \lambda_1, \lambda_2, \theta_{01}, \theta_{02}, \theta_{11}, \theta_{12}, \mu_2, C, \alpha_1, \alpha_2)$, where the platform's decision variables are the matching thresholds $(\mu_{11}^*, \mu_{12}^*)$.

From an economic perspective, the platform earns revenue from completed trips.
Assume that the fixed fee charged per trip is $p_f$ and that the price charged per unit of time is $p_s$.
Expected revenue by time $t$ can then be written as
\begin{equation}
\mathbb{E}\left[\sum_{i=1}^2 \left(p_f D_{2i}(t) + p_s \int_0^t Z_{2i}(s) ds\right)\right]. \label{eq:revenue}
\end{equation}
The prices and revenue-sharing arrangements between the platform and drivers are treated as exogenous.
Our focus is on how the choice of thresholds $(\mu^*_{11}, \mu^*_{12})$ affects system performance metrics such as throughput, queue lengths, and service quality.
In this sense, our analytical results can be applied with any revenue allocation between the platform and the drivers.

\section{Performance Approximation: The Fluid Model}\label{sec:Fluid_Model}
The stochastic model in Section~\ref{sec:Stochastic_Model} captures the microscopic dynamics of the system but is analytically intractable, particularly in large-scale settings.
To obtain a tractable characterization, we develop a deterministic fluid model that approximates the stochastic system in a many-server, heavy-traffic regime.
The fluid model reduces the analysis to a system of ordinary differential equations and admits a unique steady-state equilibrium that can be computed in closed form.

Because both customer classes draw from the same pool of idle drivers, any matching policy for one class affects the pickup efficiency of the other.
The fluid equilibrium makes this coupling explicit: the steady-state queue lengths, pickup rates, and driver allocations are jointly determined by both thresholds $(\mu_{11}^*, \mu_{12}^*)$.
This structure provides the foundation for the optimization analysis in Section~\ref{sec:matching_policies}.

%This section develops this strategic tool by deriving the fluid model as the many-server, 
%heavy-traffic limit of the stochastic system. We establish this approximation within a 
%formal scaling regime. Consider a sequence of stochastic systems indexed by the total 
%number of drivers, $N$. In the $N$-th system, passenger arrival rates are scaled 
%proportionally to the system size, i.e., $\lambda_i^N = N\lambda_i$, while all other 
%service and patience rates ($\mu_2, \theta_{0i}, \theta_{1i}$) remain of order one. 
%Crucially, to prevent the pickup process from becoming instantaneous in the limit—a 
%scenario that would obscure the vital role of spatial friction—we scale the spatial 
%matching constant as $C^N = C/N^{\alpha_1 + \alpha_2}$. This ensures that the pickup 
%rates, $\mu_{1i}(t)$, remain of constant order. Under this regime, we rigorously prove 
%in Appendix EC.2 that the scaled stochastic processes (e.g., $Q_i(t)/N, Z_{ji}(t)/N$) 
%converge in probability to their deterministic fluid counterparts.

%Revised paragraph on scaling regime:
%\paragraph{Scaling Regime (CN Scaling).}
To develop this approximation, we analyze a sequence of stochastic systems indexed by the total number of drivers $N$.
In the $N$-th system, passenger arrival rates grow proportionally with system size, $\lambda_i^N = N\lambda_i$, while all service and patience rates ($\mu_2, \theta_{0i}, \theta_{1i}$) remain of order one.

A subtlety arises in the spatial pickup process: as $N$ increases, the average distance between idle drivers and waiting passengers would shrink roughly as $N^{-(\alpha_1+\alpha_2)}$ under standard spatial intensity assumptions.
Without proper rescaling, the pickup process would become instantaneous in the limit, eliminating the critical role of endogenous pickup times that shape system performance.
To preserve this nontrivial spatial effect, we therefore scale the matching constant as
\begin{equation*}
  C^N = \frac{C}{N^{\alpha_1+\alpha_2}},
\end{equation*}
which keeps the effective pickup rates $\mu_{1i}(t)$ of constant order. 

Under this regime, the scaled stochastic processes (e.g., $Q_i(t)/N$, $Z_{ji}(t)/N$) converge in probability to their deterministic fluid counterparts, as formally established in Appendix~\ref{sec:EC of fluid model convergence} for the symmetric-exponent case $\alpha_1 = \alpha_2$.\footnote{The convergence proof in Appendix~\ref{sec:EC of fluid model convergence} is carried out under $\alpha_1 = \alpha_2$, which is the empirically relevant case \citep{WangZhangHu2024}. All numerical experiments in this paper use $\alpha_1 = \alpha_2 = 0.5$.}
The CN scaling thus provides a consistent bridge between the stochastic dynamics and the fluid limit: as the market grows, the expected pickup distance shrinks at rate $N^{-(\alpha_1+\alpha_2)}$, and the rescaling of $C$ compensates for this contraction so that endogenous pickup times remain of first-order importance in the limiting dynamics.

\subsection{Fluid Model Formulation}
\label{sec:fluid_model_formulation}
The fluid model describes the evolution of the proportions of passengers and drivers in various states.
Let $q_i(t)$ be the fluid-scaled number of waiting class-$i$ customers, $i \in \{1, 2\}$, and let $z_0(t)$, $z_{1i}(t)$, and $z_{2i}(t)$ be the proportions of drivers who are idle, assigned to class $i$, and busy with class $i$, respectively.
As these are proportions of the total driver pool, they must satisfy the conservation law for all $t \ge 0$:
\begin{equation}
z_0(t) + z_{11}(t) + z_{12}(t) + z_{21}(t) + z_{22}(t) = 1.
\label{eq:fluid_conservation}
\end{equation}

The dynamics of the system are captured by a set of ordinary differential equations (ODEs).
The core idea is to replace stochastic processes with their mean rates of change.
Let $\pi_i(t)$ denote the instantaneous matching rate for class-$i$ customers in the fluid limit.
The cumulative number of class-$i$ matches up to time $t$, denoted by $m_i(t)$, is given by $m_i(t) = \int_0^t \pi_i(s) ds$.
The evolution of the system state is governed by the following equations, which describe how the number of waiting customers, idle drivers, and drivers in various states change over time:
\begin{align}
q'_1(t) &= \lambda_1 - \theta_{01}q_1(t) - \pi_1(t), \label{eq:fluid_q1}\\
q'_2(t) &= \lambda_2 - \theta_{02}q_2(t) - \pi_2(t), \label{eq:fluid_q2}\\
z'_0(t) &= \theta_{11}z_{11}(t) + \theta_{12}z_{12}(t) + \mu_2(z_{21}(t) + z_{22}(t)) - (\pi_1(t) + \pi_2(t)), \label{eq:fluid_z0}\\
z'_{11}(t) &= \pi_1(t) - (\mu_{11}(t)+\theta_{11})z_{11}(t), \label{eq:fluid_z11}\\
z'_{12}(t) &= \pi_2(t) - (\mu_{12}(t)+\theta_{12})z_{12}(t), \label{eq:fluid_z12}\\
z'_{21}(t) &= \mu_{11}(t)z_{11}(t) - \mu_2 z_{21}(t), \label{eq:fluid_z21}\\
z'_{22}(t) &= \mu_{12}(t)z_{12}(t) - \mu_2 z_{22}(t). \label{eq:fluid_z22}
\end{align}
Here, the pickup rates $\mu_{11}(t)$ and $\mu_{12}(t)$ are endogenously determined by the system state through the spatial matching functions:
\begin{align}
\mu_{11}(t) &= C q_1(t)^{\alpha_1} z_0(t)^{\alpha_2}, \label{eq:fluid_mu11}\\
\mu_{12}(t) &= C q_2(t)^{\alpha_1} z_0(t)^{\alpha_2}. \label{eq:fluid_mu12}
\end{align}

The matching rates, $\pi_1(t)$ and $\pi_2(t)$, are control variables determined by the platform's policy.
We analyze a threshold-based policy, where the platform sets target pickup rates, $\mu_{11}^*$ and $\mu_{12}^*$, for each customer class.
Matching for a given class occurs only to prevent its pickup rate from exceeding the corresponding threshold.
This policy can be stated formally as:
\begin{gather}
\mu_{1i}(t) \leq \mu_{1i}^*, \quad \forall t \ge 0, \quad i \in \{1, 2\}, \label{eq:policy_cond1} \\
\int_0^t \mathbf{1}_{\{\mu_{1i}(s) < \mu_{1i}^*\}} d\pi_i(s) = 0, \quad \forall t \ge 0, \quad i \in \{1, 2\}. \label{eq:policy_cond2}
\end{gather}
Condition \eqref{eq:policy_cond1} is a state constraint ensuring the pickup rate never surpasses the threshold.
Condition \eqref{eq:policy_cond2} is a complementarity condition, stating that the matching rate $\pi_i(t)$ can be positive only when the pickup rate $\mu_{1i}(t)$ is precisely at its threshold $\mu_{1i}^*$.
This policy is intuitive: the platform defers matching until the pickup rate, which is inversely related to the expected pickup time, reaches its class-specific target.

To provide an alternative view of the system dynamics, we express the state variables in integral form.
This representation clarifies the link between stochastic and deterministic dynamics.
Let $d_{1i}(t)$ and $r_{1i}(t)$ be the cumulative number of successful pickups and post-matching cancellations for class-$i$ customers up to time $t$, respectively.
These quantities depend on the entire history of the matching process $m_i(t)$.

For drivers who are in the ``assigned to class $i$'' state at time $t=0$, with initial quantity $z_{1i}(0)$, the number who have successfully picked up their passenger by time $t$ is given by $\frac{\mu_{1i}(0)}{\mu_{1i}(0) + \theta_{1i}} z_{1i}(0)(1 - e^{-(\mu_{1i}(0) + \theta_{1i})t})$.
For drivers matched at a later time $s$ in the interval $0 < s \leq t$, the probability of a successful pickup is $\frac{\mu_{1i}(s)}{\mu_{1i}(s) + \theta_{1i}}$.
Integrating over all past matches, we obtain the complete expression for the cumulative successful pickups $d_{1i}(t)$ and cancellations $r_{1i}(t)$:

\begin{align}
d_{1i}(t) &= \frac{\mu_{1i}(0)}{\mu_{1i}(0) + \theta_{1i}} z_{1i}(0)(1 - e^{-(\mu_{1i}(0) + \theta_{1i})t}) \nonumber \\
&\quad + \int_0^t \frac{\mu_{1i}(s)}{\mu_{1i}(s) + \theta_{1i}}(1 - e^{-(\mu_{1i}(s) + \theta_{1i})(t-s)}) \pi_i(s) ds, \label{eq:integral_di} \\
r_{1i}(t) &= \frac{\theta_{1i}}{\mu_{1i}(0) + \theta_{1i}} z_{1i}(0)(1 - e^{-(\mu_{1i}(0) + \theta_{1i})t}) \nonumber \\
&\quad + \int_0^t \frac{\theta_{1i}}{\mu_{1i}(s) + \theta_{1i}}(1 - e^{-(\mu_{1i}(s) + \theta_{1i})(t-s)}) \pi_i(s) ds. \label{eq:integral_ri}
\end{align}

Here, we assume for simplicity that the pickup rate for the initially assigned drivers $z_{1i}(0)$ is determined by the initial state $\mu_{1i}(0)$.
Under our threshold policy, for any $s$ where matching occurs ($\pi_i(s)>0$), we have $\mu_{1i}(s) = \mu_{1i}^*$.

With these definitions, the system dynamics can be expressed as a set of integral equations:
\begin{align*}
q_i(t) &= q_i(0) + \lambda_i t - \theta_{0i} \int_0^t q_i(s) ds - m_i(t), \\
z_{1i}(t) &= z_{1i}(0) + m_i(t) - d_{1i}(t) - r_{1i}(t), \\
z_{2i}(t) &= z_{2i}(0) + d_{1i}(t) - \mu_2 \int_0^t z_{2i}(s) ds.
\end{align*}

By taking the time derivative of these integral equations (and applying the Leibniz rule to equations \eqref{eq:integral_di} and \eqref{eq:integral_ri}), we recover the system of ordinary differential equations (ODEs) \eqref{eq:fluid_q1}-\eqref{eq:fluid_z22}.
For instance, $z'_{1i}(t) = m'_i(t) - d'_{1i}(t) - r'_{1i}(t) = \pi_i(t) - (\mu_{1i}(t)+\theta_{1i})z_{1i}(t)$.
This demonstrates that our ODE model is a direct and more tractable representation of these underlying integral dynamics.

With all components defined, we can now formally state the fluid model.
\begin{definition}[Fluid Process]
For a given set of parameters and decision thresholds $(\mu_{11}^*, \mu_{12}^*)$, a process $x(t) = (q_1(t), q_2(t), z_0(t), z_{11}(t), z_{12}(t), z_{21}(t), z_{22}(t))$ is a fluid solution if it satisfies the conservation law \eqref{eq:fluid_conservation} and the dynamic equations \eqref{eq:fluid_q1}-\eqref{eq:policy_cond2}.
\end{definition}

\begin{proposition} \label{prop:fluid_model_existence_uniqueness}
The fluid model for the two-class customer system, as defined by equations \eqref{eq:fluid_conservation}-\eqref{eq:policy_cond2} with the threshold matching policy where matching for class $i \in \{1,2\}$ occurs if its potential matching intensity $\mu_{1i}(t) = C q_i(t)^{\alpha_1} z_0(t)^{\alpha_2}$ meets or exceeds a threshold $\mu_{1i}^*$, has a unique solution trajectory $(q_i(t), z_0(t), z_{1i}(t), z_{2i}(t))$ for $t \ge 0$ given initial conditions.
\end{proposition}
A rigorous proof of existence and uniqueness of the fluid model solution is provided in Appendix~\ref{sec:EC of fluid model convergence}.
The argument proceeds by first establishing that the system of ODEs under the threshold policy forms a well-posed initial value problem with locally Lipschitz right-hand sides, ensuring local existence and uniqueness via the Picard–Lindelöf theorem~\citep{coddington1955theory}.
Global existence and boundedness follow from the conservation law and non-negativity constraints.
Uniqueness of the solution trajectory is further established by showing that the mapping defined by the integral form of the dynamics is a contraction under suitable norms, exploiting the monotonicity and continuity properties of the system.

\subsection{Equilibrium Analysis and Approximation}
\label{sec:fluid_model_equilibrium}
A key advantage of the fluid model is that it often converges to a unique, stable equilibrium.
This equilibrium represents the long-run average behavior of the system and can be found by solving a system of algebraic equations, which is far more tractable than analyzing the stochastic model's steady-state distribution.

At equilibrium, the system is in steady state, so all time derivatives of the state variables are zero: $q'_i(t)=0, z'_{ji}(t)=0$.
Setting the left-hand side of equations \eqref{eq:fluid_q1}-\eqref{eq:fluid_z22} to zero yields a set of flow-balance equations.
Before stating the main convergence result, we impose feasibility requirements for a non-degenerate equilibrium: the pre-match abandonment rates must be strictly positive ($\theta_{0i} > 0$) and the thresholds must be attainable.
Because $\theta_{0i}>0$, the queue length is bounded in the absence of matching by $q_i(t)\le \lambda_i/\theta_{0i}$, and since $z_0(t)\le 1$, the pickup rate $\mu_{1i}(t)=C q_i(t)^{\alpha_1} z_0(t)^{\alpha_2}$ has a natural upper bound $\mu_{1i}(t)\le C(\lambda_i/\theta_{0i})^{\alpha_1}$.
Thus we restrict attention to policies with $\mu_{1i}^*\in (0,\, C(\lambda_i/\theta_{0i})^{\alpha_1}]$; setting a larger threshold would preclude matching for class $i$.

%Second, we impose a mild condition on the system's time scales: 
%$\theta_{1i} > \mu_2$ for $i=1, 2$. This assumption reflects 
%the typical real-world scenario in which the average time a 
%passenger is willing to wait for pickup after matching is shorter 
%than the average duration of the trip itself. This condition ensures 
%a stable flow between the "assigned" and "busy" driver states and 
%prevents an unrealistic accumulation of drivers in long trips relative 
%to their pickup phase. The following theorem characterizes the unique 
%equilibrium that exists under these well-posed conditions.

\begin{theorem}[Convergence to Equilibrium] \label{thm:equilibrium}
Consider the fluid model where matching for class $i \in \{1,2\}$ occurs
if its potential matching intensity $\mu_{1i}(t) = C q_i(t)^{\alpha_1} z_0(t)^{\alpha_2}$ meets
or exceeds a positive threshold $\mu_{1i}^* > 0$.
If system parameters allow for a physical equilibrium where matching is active
for both classes, that is, the derived matching rates $\pi_i^*$ are non-negative and $z_0^* > 0$, 
and the conditions $\mu_{1i}^* \in \left(0, C\left(\frac{\lambda_i}{\theta_{0i}}\right)^{\alpha_1}\right]$ 
hold for $i=1,2$,
then there exists a unique non-negative solution  
$(q_1^*, q_2^*, z_0^*, z_{11}^*, z_{12}^*, z_{21}^*, z_{22}^*)$ 
with associated non-negative matching rates $(\pi_1^*, \pi_2^*)$ to the equilibrium equations:
    \begin{align}
     \lambda_1 &= q_1^*\theta_{01} + z_{11}^*\theta_{11} + z_{21}^*\mu_2, \label{eq:bal_q1}\\
     \lambda_2 &= q_2^*\theta_{02} + z_{12}^*\theta_{12} + z_{22}^*\mu_2, \label{eq:bal_q2}\\
     z_{11}^*\mu_{11}^* &= z_{21}^*\mu_2, \label{eq:bal_z11}\\
     z_{12}^*\mu_{12}^* &= z_{22}^*\mu_2, \label{eq:bal_z12}\\
     1 &= z_0^* + z_{11}^* + z_{12}^* + z_{21}^* + z_{22}^*, \label{eq:bal_cons}\\
     \mu_{11}^* &= C (q_1^*)^{\alpha_1}(z_0^*)^{\alpha_2}, \label{eq:bal_mu11}\\
     \mu_{12}^* &= C (q_2^*)^{\alpha_1}(z_0^*)^{\alpha_2}. \label{eq:bal_mu12}
     \end{align}
This solution is the unique physical equilibrium point of the fluid model 
when both matching thresholds $\mu_{1i}^*$ are active. In addition,
the fluid process $(q_i(t), z_0(t), z_{1i}(t), z_{2i}(t))$ for $i=1,2$, 
converges to this equilibrium point exponentially fast as $t \to \infty$.
\end{theorem}

The intuition behind these equilibrium equations is straightforward. Equations 
\eqref{eq:bal_q1} and \eqref{eq:bal_q2} represent the passenger flow balance for each class: 
the arrival rate must equal the total rate of departure through abandonment before matching 
($q_i^*\theta_{0i}$), cancellation after matching ($z_{1i}^*\theta_{1i}$),
and successful trip completion ($z_{2i}^*\mu_2$). Equations \eqref{eq:bal_z11}
and \eqref{eq:bal_z12} represent the driver flow balance between the ``assigned'' and
``busy'' states: the rate at which drivers become busy, i.e., successful pickups ($z_{1i}^*\mu_{1i}^*$),
must equal the rate at which they leave the busy state, i.e., trip completions ($z_{2i}^*\mu_2$). 
Equation \eqref{eq:bal_cons} is the driver conservation law at steady state. Finally, equations 
\eqref{eq:bal_mu11} and \eqref{eq:bal_mu12} enforce the matching policy at equilibrium. 
The proof of this theorem is detailed in Appendix~\ref{sec:EC of fluid model convergence}.

\subsection{Numerical Validation and Strategic Insights}
\label{sec:numerical_validation_and_insights}

This section evaluates the accuracy of the fluid approximation and illustrates how the model reveals
operational implications for platforms managing heterogeneous customer classes. We assess the
approximation along two dimensions (process-level convergence and steady-state accuracy) and use
the equilibrium structure to examine the effects of threshold-based differentiation.

We consider a system with $N=1000$ drivers and the baseline parameters summarized in
Table~\ref{tab:baseline_parameters}. Service differentiation is implemented by varying the matching
thresholds $(\mu_{11}^*, \mu_{12}^*)$ across three representative policies:
\begin{equation*}
  \text{A}: (6,6),\qquad
  \text{B}: (4,8),\qquad
  \text{C}: (8,4),
\end{equation*}
where smaller thresholds correspond to higher service standards.

\begin{table}[htbp]
\centering
\caption{Baseline Parameters for Numerical Experiments}
\label{tab:baseline_parameters}
\begin{tabular}{lcc}
\toprule
Parameter & Description & Value \\
\midrule
$\lambda_1, \lambda_2$ & Arrival rates (Class 1, Class 2) & 2.0,\ 3.0 \\
$\theta_{01}, \theta_{02}$ & Pre-match abandonment rates & 2.0,\ 1.0 \\
$\theta_{11}, \theta_{12}$ & Post-match cancellation rates & 3.0,\ 4.0 \\
$\mu_2$ & Trip completion rate & 10.0 \\
$C$ & Spatial matching constant & 50.0 \\
$\alpha_1, \alpha_2$ & Spatial exponents & 0.5,\ 0.5 \\
\bottomrule
\end{tabular}
\end{table}

%\paragraph{Process-Level Validation.}
Figure~\ref{fig:fluid_vs_stochastic} plots a sample path of the scaled stochastic process and the corresponding fluid trajectory under Scenario~A.
The two dynamics nearly coincide, including their transient behavior and convergence rate.
Similar agreement holds under Scenarios~B and~C, as well as at higher load levels, indicating that the CN-scaling fluid model captures the essential system evolution beyond lightly loaded regimes.

\begin{figure}[htbp]
    \centering
    \includegraphics[width=0.9\textwidth]{figs/fluid_vs_stochastic_final_unified_legend.png}
    \caption{Fluid and Stochastic Trajectories under Scenario~A ($N=1000$).}
    \label{fig:fluid_vs_stochastic}
\end{figure}

%\paragraph{Steady-State Validation.}
Table~\ref{tab:policy_and_load_comparison} compares steady-state simulations for increasing system sizes with the fluid equilibrium from Theorem~\ref{thm:equilibrium}.
Two patterns consistently emerge: (i) scaled stochastic averages converge numerically to the fluid predictions, and (ii) the variance diminishes as $N$ grows.
For instance, in Scenario~A (high load) the simulated ratios $E[Q_1]/N=0.0710$ and $E[Q_2]/N=0.0712$ at $N=1000$ differ from the fluid values by less than $1\%$.
Across all scenarios and load conditions, queue-length discrepancies remain within $6\%$, supporting the model's reliability for strategic analysis.

\begin{table}[htbp]
\centering
\caption{Comparison of Fluid Approximation and Stochastic System in Steady State under Different Policies and Loads.}
\label{tab:policy_and_load_comparison}
 \begin{threeparttable}
\resizebox{\textwidth}{!}{%
\begin{tabular}{l|l|l|r@{\,$\pm$\,}l|r@{\,$\pm$\,}l|r@{\,$\pm$\,}l|r@{\,$\pm$\,}l|r@{\,$\pm$\,}l|r@{\,$\pm$\,}l|r@{\,$\pm$\,}l}
\hline
\textbf{Scenario} & \textbf{Load} & \textbf{N} & \multicolumn{2}{c|}{\textbf{$E[Q_1]/N$}} & \multicolumn{2}{c|}{\textbf{$E[Q_2]/N$}} & \multicolumn{2}{c|}{\textbf{$E[Z_0]/N$}} & \multicolumn{2}{c|}{\textbf{$E[Z_{11}]/N$}} & \multicolumn{2}{c|}{\textbf{$E[Z_{12}]/N$}} & \multicolumn{2}{c|}{\textbf{$E[Z_{21}]/N$}} & \multicolumn{2}{c}{\textbf{$E[Z_{22}]/N$}} \\ \hline
\multirow{8}{*}{A (Symmetric)} & \multirow{4}{*}{$\lambda \times 0.5$} & 100 & 0.0100 & 0.0000 & 0.0102 & 0.0000 & 0.7971 & 0.0010 & 0.0539 & 0.0004 & 0.0736 & 0.0006 & 0.0317 & 0.0004 & 0.0437 & 0.0005 \\ 
 & & 500 & 0.0168 & 0.0000 & 0.0169 & 0.0000 & 0.8012 & 0.0005 & 0.0513 & 0.0003 & 0.0729 & 0.0003 & 0.0308 & 0.0002 & 0.0437 & 0.0003 \\ 
 & & 1000 & 0.0174 & 0.0000 & 0.0175 & 0.0000 & 0.8015 & 0.0004 & 0.0510 & 0.0002 & 0.0732 & 0.0002 & 0.0306 & 0.0001 & 0.0438 & 0.0002 \\ 
 & & \textbf{Fluid} & \multicolumn{2}{c|}{\textbf{0.0180}} & \multicolumn{2}{c|}{\textbf{0.0180}} & \multicolumn{2}{c|}{\textbf{0.8004}} & \multicolumn{2}{c|}{\textbf{0.0516}} & \multicolumn{2}{c|}{\textbf{0.0732}} & \multicolumn{2}{c|}{\textbf{0.0309}} & \multicolumn{2}{c}{\textbf{0.0439}} \\ 
\cline{2-17}
 & \multirow{4}{*}{$\lambda \times 2.0$} & 100 & 0.0718 & 0.0011 & 0.0729 & 0.0011 & 0.2038 & 0.0026 & 0.2082 & 0.0011 & 0.2894 & 0.0013 & 0.1242 & 0.0008 & 0.1744 & 0.0012 \\ 
 & & 500 & 0.0705 & 0.0004 & 0.0707 & 0.0004 & 0.2075 & 0.0011 & 0.2051 & 0.0004 & 0.2904 & 0.0007 & 0.1225 & 0.0004 & 0.1745 & 0.0004 \\ 
 & & 1000 & 0.0710 & 0.0002 & 0.0712 & 0.0002 & 0.2066 & 0.0005 & 0.2058 & 0.0003 & 0.2907 & 0.0003 & 0.1228 & 0.0002 & 0.1742 & 0.0003 \\ 
 & & \textbf{Fluid} & \multicolumn{2}{c|}{\textbf{0.0715}} & \multicolumn{2}{c|}{\textbf{0.0715}} & \multicolumn{2}{c|}{\textbf{0.2013}} & \multicolumn{2}{c|}{\textbf{0.2063}} & \multicolumn{2}{c|}{\textbf{0.2928}} & \multicolumn{2}{c|}{\textbf{0.1238}} & \multicolumn{2}{c}{\textbf{0.1757}} \\ 
\hline
\multirow{8}{*}{B (Prioritize VIPs)} & \multirow{4}{*}{$\lambda \times 0.5$} & 100 & 0.0000 & 0.0000 & 0.0287 & 0.0001 & 0.7929 & 0.0016 & 0.0716 & 0.0009 & 0.0599 & 0.0007 & 0.0282 & 0.0004 & 0.0474 & 0.0006 \\ 
 & & 500 & 0.0071 & 0.0000 & 0.0311 & 0.0000 & 0.7969 & 0.0006 & 0.0689 & 0.0004 & 0.0593 & 0.0003 & 0.0274 & 0.0002 & 0.0475 & 0.0002 \\ 
 & & 1000 & 0.0076 & 0.0000 & 0.0316 & 0.0000 & 0.7966 & 0.0006 & 0.0690 & 0.0003 & 0.0595 & 0.0002 & 0.0276 & 0.0001 & 0.0473 & 0.0002 \\ 
 & & \textbf{Fluid} & \multicolumn{2}{c|}{\textbf{0.0080}} & \multicolumn{2}{c|}{\textbf{0.0322}} & \multicolumn{2}{c|}{\textbf{0.7955}} & \multicolumn{2}{c|}{\textbf{0.0691}} & \multicolumn{2}{c|}{\textbf{0.0598}} & \multicolumn{2}{c|}{\textbf{0.0277}} & \multicolumn{2}{c}{\textbf{0.0479}} \\ 
\cline{2-17}
 & \multirow{4}{*}{$\lambda \times 2.0$} & 100 & 0.0322 & 0.0006 & 0.1449 & 0.0023 & 0.1863 & 0.0027 & 0.2761 & 0.0019 & 0.2371 & 0.0010 & 0.1118 & 0.0008 & 0.1887 & 0.0009 \\ 
 & & 500 & 0.0338 & 0.0002 & 0.1383 & 0.0007 & 0.1908 & 0.0010 & 0.2745 & 0.0005 & 0.2365 & 0.0003 & 0.1093 & 0.0004 & 0.1888 & 0.0004 \\ 
 & & 1000 & 0.0341 & 0.0001 & 0.1382 & 0.0004 & 0.1901 & 0.0006 & 0.2746 & 0.0004 & 0.2374 & 0.0003 & 0.1090 & 0.0003 & 0.1890 & 0.0003 \\ 
 & & \textbf{Fluid} & \multicolumn{2}{c|}{\textbf{0.0347}} & \multicolumn{2}{c|}{\textbf{0.1386}} & \multicolumn{2}{c|}{\textbf{0.1847}} & \multicolumn{2}{c|}{\textbf{0.2758}} & \multicolumn{2}{c|}{\textbf{0.2384}} & \multicolumn{2}{c|}{\textbf{0.1103}} & \multicolumn{2}{c}{\textbf{0.1908}} \\ 
\hline
\multirow{8}{*}{C (Prioritize Regulars)} & \multirow{4}{*}{$\lambda \times 0.5$} & 100 & 0.0276 & 0.0001 & 0.0000 & 0.0000 & 0.7977 & 0.0018 & 0.0406 & 0.0004 & 0.0923 & 0.0010 & 0.0317 & 0.0006 & 0.0376 & 0.0006 \\ 
 & & 500 & 0.0308 & 0.0000 & 0.0070 & 0.0000 & 0.8000 & 0.0004 & 0.0393 & 0.0002 & 0.0921 & 0.0003 & 0.0316 & 0.0001 & 0.0370 & 0.0002 \\ 
 & & 1000 & 0.0314 & 0.0000 & 0.0075 & 0.0000 & 0.7996 & 0.0005 & 0.0395 & 0.0001 & 0.0925 & 0.0003 & 0.0317 & 0.0001 & 0.0368 & 0.0001 \\ 
 & & \textbf{Fluid} & \multicolumn{2}{c|}{\textbf{0.0320}} & \multicolumn{2}{c|}{\textbf{0.0080}} & \multicolumn{2}{c|}{\textbf{0.7988}} & \multicolumn{2}{c|}{\textbf{0.0396}} & \multicolumn{2}{c|}{\textbf{0.0927}} & \multicolumn{2}{c|}{\textbf{0.0317}} & \multicolumn{2}{c}{\textbf{0.0371}} \\ 
\cline{2-17}
 & \multirow{4}{*}{$\lambda \times 2.0$} & 100 & 0.1324 & 0.0018 & 0.0304 & 0.0005 & 0.1969 & 0.0023 & 0.1574 & 0.0008 & 0.3721 & 0.0017 & 0.1255 & 0.0009 & 0.1481 & 0.0012 \\ 
 & & 500 & 0.1301 & 0.0006 & 0.0320 & 0.0002 & 0.2010 & 0.0009 & 0.1564 & 0.0004 & 0.3694 & 0.0007 & 0.1251 & 0.0004 & 0.1481 & 0.0004 \\ 
 & & 1000 & 0.1297 & 0.0005 & 0.0321 & 0.0001 & 0.2015 & 0.0007 & 0.1565 & 0.0003 & 0.3695 & 0.0005 & 0.1252 & 0.0003 & 0.1473 & 0.0002 \\ 
 & & \textbf{Fluid} & \multicolumn{2}{c|}{\textbf{0.1305}} & \multicolumn{2}{c|}{\textbf{0.0326}} & \multicolumn{2}{c|}{\textbf{0.1962}} & \multicolumn{2}{c|}{\textbf{0.1581}} & \multicolumn{2}{c|}{\textbf{0.3709}} & \multicolumn{2}{c|}{\textbf{0.1265}} & \multicolumn{2}{c}{\textbf{0.1484}} \\ 
\hline
\end{tabular}%
}
\begin{tablenotes}[flushleft]
  \footnotesize
  \item \textit{Notes.} Base parameters: $\lambda_1=1.0$, $\lambda_2=1.5$, $\theta_{01}=2.0$, $\theta_{02}=1.0$, 
  $\theta_{11}=3.0$, $\theta_{12}=4.0$, $C=50.0$, $\alpha_1=\alpha_2=0.5$, $\mu_2=10.0$.
  \item Matching thresholds $(\mu_{11}^*, \mu_{12}^*)$ vary by scenario: A (6,6), B (4,8), C (8,4).
\end{tablenotes}
\end{threeparttable}
\end{table}

%\paragraph{Strategic Insights.}
The numerical experiments reveal how the shared pool of idle drivers mediates cross-class externalities.
As indicated by the equilibrium matching conditions \eqref{eq:bal_mu11}--\eqref{eq:bal_mu12}, both customer classes rely on the same endogenous idle-driver level $z_0^*$; consequently, raising the service standard for one segment inevitably constrains the effective matching rate for the other.
This trade-off is clearly visible under high load when shifting from the symmetric benchmark (Scenario~A) to a policy prioritizing VIPs (Scenario~B).
In this case, the VIP queue is reduced by more than half ($0.0715\to 0.0347$, a $51\%$ reduction), yet this improvement comes at a significant cost to the regular queue ($0.0715\to 0.1386$, a $94\%$ increase), demonstrating the cross-class externalities inherent in shared-capacity systems.

Beyond simple trade-offs, the system exhibits distinct structural behaviors under asymmetric prioritization.
For instance, under Scenario~C at high load, the aggressive threshold set for regular customers ($\mu_{12}^*=4$) reduces the idle-driver pool ($z_0^*=0.1962$) while substantially increasing the VIP queue ($q_1^*=0.1305$, compared to $0.0715$ in Scenario~A).
This illustrates how prioritizing one class can significantly impact waiting times for the other class.
The convergence of stochastic simulations toward these fluid predictions confirms that these trade-offs are intrinsic properties of the system dynamics, not artifacts of the approximation.

From a strategic perspective, these findings suggest that threshold-based policies $(\mu_{11}^*,\mu_{12}^*)$ offer greater flexibility than binary priority rules, providing a continuous mechanism to regulate the allocation of idle supply.
Small adjustments in thresholds translate into predictable shifts in equilibrium queue lengths and pickup rates, enabling platforms to quantify the trade-off between revenue and customer retention.
The fluid model thus provides both an accurate approximation for transient and steady-state performance and a tractable foundation for optimizing service levels across heterogeneous customer classes.

%\begin{remark}
%As shown in Table 1, our approximation works reasonably well, even at a scale of $n = 100$. But we also observe that the accuracy of our approximation decreases as $E(Q)$ or $E(Z_0)$ becomes small primarily because of the subsequent rounding loss. In the original stochastic process, $\mu_1(t) = \tilde{C}(Q(t))^{\alpha_1}(Z_0(t))^{\alpha_2}$ takes discrete
%\end{remark}

% =============================================================================
% SECTION 4: OPTIMAL MATCHING POLICY DESIGN
% =============================================================================
% Structure Overview:
% 4.1 Problem Formulation - Sets up the optimization problem (~0.5 page)
% 4.2 Optimality Conditions - Derives KMI via KKT conditions (~1.5 pages)
%     - Proposition: Idle-driver marginal value conservation
%     - Definition: Key Matching Indices
%     - Theorem: Stationarity characterization
% 4.3 Sensitivity Analysis - Numerical properties of Jacobian (~0.5 page)
% 4.4 Coordinated Control Algorithm - Newton-Raphson implementation (~0.5 page)
% 4.5 Strategic Insights - Managerial implications (~0.5 page)
% =============================================================================

\section{Optimal Matching Policy Design}\label{sec:matching_policies}

% =============================================================================
% OPENING PARAGRAPH: Transition from fluid model to policy design
% Purpose: Connect to previous section and motivate the optimization problem
% =============================================================================
Having established the fluid model and its steady-state properties in Section~\ref{sec:Fluid_Model}, we now use this framework to design optimal matching policies.
The fluid approximation provides a tractable analytical basis for characterizing the interplay between matching thresholds, driver supply, and passenger waiting times.
This section formulates the threshold optimization problem and derives the necessary optimality conditions.

% =============================================================================
% SECOND PARAGRAPH: Preview of contribution and three forces
% Purpose: Introduce the key trade-off and the KMI framework
% =============================================================================
The central prescriptive question is: \emph{how should the platform dynamically adjust the thresholds $(\mu_{11}^*, \mu_{12}^*)$ to optimize a given objective?}
The answer involves a careful balancing of three competing forces: \emph{(i)} the benefit of reducing customer abandonment by matching earlier, \emph{(ii)} the opportunity cost of committing a driver to one class, and \emph{(iii)} the externality imposed on the other class through the shared idle-driver pool.
These trade-offs are encoded in the \emph{Key Matching Indices} (KMI), from which we derive a coordinated control algorithm capable of maximizing platform objectives in real time.

\subsection{Problem Formulation}
\label{subsec:problem_formulation}

% =============================================================================
% PARAGRAPH: Decision variables and objective
% Purpose: Formally state the optimization problem
% =============================================================================
We study the prescriptive design of matching policies by characterizing the optimal 
thresholds in the fluid limit.
To streamline the exposition, we introduce the vector notation 
$\boldsymbol{\mu} = (\mu_{11}, \mu_{12})^\top \in \mathbb{R}_{++}^2$ for the 
threshold pair, where each component $\mu_{1i}$ corresponds to the threshold 
$\mu_{1i}^*$ introduced in Section~\ref{sec:Fluid_Model}.
Let $\mathbf{x}^* = (z_0^*, z_{11}^*, z_{12}^*, z_{21}^*, z_{22}^*, q_1^*, q_2^*)$ 
denote the unique steady state induced by a given $\boldsymbol{\mu}$, as characterized 
in Theorem~\ref{thm:equilibrium}.
The platform seeks to maximize a weighted throughput objective:
\begin{equation}
\max_{\boldsymbol{\mu} > 0} 
\quad 
J(\boldsymbol{\mu})
    := 
    w_1 z_{21}^*(\boldsymbol{\mu})
    + 
    w_2 z_{22}^*(\boldsymbol{\mu}),
\label{eq:opt_problem}
\end{equation}
where $w_i$ denotes the economic weight of class $i$.
Setting $w_1 = w_2 = 1$ yields a \emph{social welfare} objective that maximizes 
the total rate of successfully completed trips.
Alternatively, setting $w_i = p_i$, where $p_i$ is the price charged to class $i$,
yields a \emph{platform revenue} objective.
Since the trip completion rate for class $i$ is $\mu_2 z_{2i}^*$, the objective $\sum_i w_i z_{2i}^*$ is proportional to weighted throughput; setting $w_i = p_i$ thus weights each class's contribution by its per-trip revenue.

% =============================================================================
% PARAGRAPH: Implicit dependence and complexity
% Purpose: Explain why the problem is nontrivial
% =============================================================================
The challenge in solving~\eqref{eq:opt_problem} stems from the implicit dependence 
of $\mathbf{x}^*$ on $\boldsymbol{\mu}$ through the nonlinear equilibrium 
conditions~\eqref{eq:bal_q1}--\eqref{eq:bal_mu12}.
Standard gradient-based methods require evaluating the sensitivity 
$\nabla_{\boldsymbol{\mu}} \mathbf{x}^*$, which is not available in closed form.
Our approach circumvents this difficulty by deriving compact
\emph{optimality indicators}, the Key Matching Indices, that encapsulate the first-order conditions
in terms of observable steady-state quantities.

\subsection{Optimality Conditions and Key Matching Indices}
\label{subsec:optimality_conditions}

% =============================================================================
% PARAGRAPH: Overview of the derivation approach
% Purpose: Explain the methodology before presenting results
% =============================================================================
We derive the first-order necessary conditions for problem~\eqref{eq:opt_problem} using the KKT framework applied to the fluid equilibrium.
The derivation, detailed in Appendix~\ref{sec:ec_optimality}, constructs a Lagrangian incorporating three classes of constraints: the passenger flow balance for each class, the driver flow balance between pickup and trip states, and the overall driver conservation law.
The Lagrange multipliers associated with these constraints admit natural economic interpretations as shadow prices, and their manipulation yields the optimality conditions stated below.

% =============================================================================
% PROPOSITION: Marginal value conservation
% Purpose: State the first key structural result
% =============================================================================
The first result establishes a necessary aggregate condition that any optimal solution must satisfy.

\begin{proposition}[Idle-Driver Marginal Value Conservation]
\label{prop:first_level}
Any interior optimal solution $\boldsymbol{\mu}$ to problem~\eqref{eq:opt_problem} must satisfy
\begin{equation}
\sum_{i=1}^2 
\frac{\alpha_2 (z_{1i}^*/z_0^*)}{s_i} = 1,
\label{eq:first_level}
\end{equation}
where 
\begin{equation*}
  s_i
  \triangleq
  1 -
  \alpha_1 \frac{z_{1i}^* \theta_{1i}}{q_i^* \theta_{0i}}
\end{equation*}
captures the elasticity of the queue length with respect to matching volume.
\end{proposition}

% =============================================================================
% PARAGRAPH: Interpretation of Proposition
% Purpose: Explain the economic meaning of the conservation law
% =============================================================================
Condition~\eqref{eq:first_level} has a natural interpretation: at optimality, the marginal value of an idle driver must be fully allocated across both matching routes.
The term $\alpha_2 (z_{1i}^*/z_0^*)/s_i$ can be understood as the shadow contribution of idle-driver supply to matching class $i$, adjusted by the queue elasticity $s_i$.
This is a spatial matching analogue of the classical equi-marginal principle: at optimality, the shadow value of an idle driver must be equalized across both matching routes.
When the conservation law is violated, the platform can improve its objective by reallocating idle-driver capacity between the two classes.

% =============================================================================
% PARAGRAPH: Transition to KMI
% Purpose: Motivate why we need more than the aggregate condition
% =============================================================================
While Proposition~\ref{prop:first_level} provides a necessary aggregate condition, it does not prescribe \emph{how} to set the individual thresholds $\mu_{11}^*$ and $\mu_{12}^*$.
To resolve this, we derive the \emph{Key Matching Indices} (KMI), which provide class-specific optimality criteria.
These indices emerge from eliminating the Lagrange multipliers between the first-order conditions and encode the trade-off between three competing effects: \emph{(i)} the benefit of reducing customer abandonment by matching earlier, \emph{(ii)} the direct opportunity cost of committing an idle driver, and \emph{(iii)} the externality imposed on the other class through the shared driver pool.

\begin{definition}[Key Matching Indices]
\label{def:kmi}
For each class $i \in \{1,2\}$, the Key Matching Index is defined as follows, with $j \neq i$ denoting the other class.
\begin{enumerate}
\item \textbf{Welfare objective} ($w_1=w_2=1$):
\begin{equation}
\zeta_i^{W}
=
\alpha_1 \frac{z_{1i}^*\theta_{1i}}{q_i^*\theta_{0i}}
+
\alpha_2\frac{z_{1i}^*}{z_0^*}
+
\frac{
\alpha_2\!\left(\frac{1}{\mu_{1i}^*} 
+ 
\alpha_1\frac{z_{1i}^*}{q_i^*\theta_{0i}}\right)
}{
\frac{z_0^*}{\mu_{1j}^*z_{1j}^*}
+
\alpha_1\frac{z_0^*}{q_j^*\theta_{0j}}
}.
\label{eq:kmi_welfare}
\end{equation}

\item \textbf{Revenue objective} ($w_i=p_i$):
\begin{equation}
\zeta_i^{R}
=
\alpha_1\frac{z_{1i}^*\theta_{1i}}{q_i^*\theta_{0i}}
+
\alpha_2\frac{z_{1i}^*}{z_0^*}
+
\alpha_2\frac{p_j}{p_i}
\frac{\mu_{1j}^*}{\mu_{1i}^*}
\frac{z_{1j}^*}{z_0^*}\,
\frac{\mu_{1i}^*s_i + \mu_2A_i}{\mu_{1j}^*s_j + \mu_2A_j},
\label{eq:kmi_revenue}
\end{equation}
where $A_i \triangleq 1 + \alpha_1 \frac{\mu_{1i}^* z_{1i}^*}{q_i^* \theta_{0i}}$ is an auxiliary elasticity factor.
\end{enumerate}
\end{definition}

% =============================================================================
% PARAGRAPH: Discussion of KMI structure
% Purpose: Provide high-level intuition for the three terms
% =============================================================================
The structure of both indices reflects the trade-offs in threshold optimization, with each term capturing one of the three competing forces:

\begin{itemize}
\item \textbf{Abandonment Reduction} (first term):
The ratio $\alpha_1 z_{1i}^*\theta_{1i}/(q_i^*\theta_{0i})$ measures the marginal benefit of reducing customer loss.
Lowering the threshold increases matching intensity, which shortens the queue and reduces both pre-match abandonment (at rate $\theta_{0i}$) and post-match cancellation (at rate $\theta_{1i}$).
This term is large when the queue is long relative to the matching flow, indicating high potential gains from more aggressive matching.

\item \textbf{Direct Opportunity Cost} (second term):
The ratio $\alpha_2 z_{1i}^*/z_0^*$ captures the cost of consuming an idle driver from the shared pool.
Each match commits a driver to the pickup-then-trip cycle, temporarily removing capacity from the system.
This term is large when picking-up drivers are numerous relative to idle drivers, signaling that the pool is already depleted.

\item \textbf{Cross-Market Externality} (third term):
The remaining term accounts for the spillover effect on the other class.
When the platform matches more aggressively for class $i$, it draws down the idle-driver pool $z_0^*$, degrading pickup efficiency for class $j$.
This externality is the key source of coupling between the two classes and distinguishes our framework from single-class models.
Examining the welfare index~\eqref{eq:kmi_welfare}, the cross-market term is large when class~$i$ matches aggressively (low threshold $\mu_{1i}^*$, so that the numerator term $1/\mu_{1i}^*$ is large) relative to the absorptive capacity of class~$j$, as measured by the idle-driver pool available per unit of class~$j$ matching flow.
Operationally, a large cross-market term signals that further threshold reduction for class~$i$ would come at a disproportionate cost to class~$j$'s service quality.
\end{itemize}

At optimality, these three effects must be balanced such that each index equals unity: $\zeta_i = 1$.

% =============================================================================
% REMARK: Welfare vs Revenue divergence
% Purpose: Highlight the structural difference between objectives
% =============================================================================
\begin{remark}[Structural Divergence Between Welfare and Revenue]
\label{remark:kmi_difference}
The revenue index $\zeta_i^R$ does not reduce to the welfare index $\zeta_i^W$ even under symmetric pricing ($p_1 = p_2 = 1$), as established in Appendix~\ref{sec:ec_kmi_difference}.
This divergence reflects an economic distinction: $\zeta_i^W$ measures opportunity costs in \emph{physical units} (throughput per hour), whereas $\zeta_i^R$ measures them in \emph{value units} (revenue per hour).
In the welfare formulation, the cross-market term penalizes the reduction in matching volume for class $j$ caused by serving class $i$.
In the revenue formulation, the Lagrange multipliers incorporate prices directly, yielding $\gamma_i \propto p_i / (\mu_{1i}^* s_i + \mu_2 A_i)$, which embeds the marginal revenue substitution rate into the index.
Consequently, even with equal prices, the revenue-optimal operating point differs from the welfare-optimal one: revenue maximization values the \emph{rate of revenue generated per driver-hour} (incorporating turnover and cancellation risk) rather than the raw \emph{volume} of completed trips.
\end{remark}

% =============================================================================
% THEOREM: Main optimality result
% Purpose: State the central characterization theorem
% =============================================================================
The following theorem establishes that the KMI provide a complete characterization of stationary points.

\begin{theorem}[Stationarity Characterization]
\label{thm:general_optimality}
A threshold vector $\boldsymbol{\mu}$ is a stationary point of problem~\eqref{eq:opt_problem} if and only if:
\begin{enumerate}
    \item the idle-driver conservation condition~\eqref{eq:first_level} holds, and
    \item both indices satisfy 
    \begin{equation*}
      \zeta_1(\boldsymbol{\mu}) = 1
      \qquad\text{and}\qquad
      \zeta_2(\boldsymbol{\mu}) = 1,
    \end{equation*}
    where $\zeta_i = \zeta_i^W$ for welfare or $\zeta_i = \zeta_i^R$ for revenue.
\end{enumerate}
The proof is provided in Appendix~\ref{sec:ec_optimality}.
\end{theorem}

\begin{remark}[On Global Optimality]
\label{remark:global_optimality}
Since the spatial matching function is non-convex, the KMI stationarity condition is necessary but not sufficient for global optimality.
Uniqueness of the fluid equilibrium (Theorem~\ref{thm:equilibrium}) and numerical evidence of unimodality (Figures~\ref{fig:welfare-3d}--\ref{fig:revenue-3d}) support the conclusion that the stationary point is the global maximizer; see Appendix~\ref{remark:sufficiency} for details.
\end{remark}

% =============================================================================
% PARAGRAPH: Interpretation and practical significance
% Purpose: Explain why this result is useful
% =============================================================================
To state the result more compactly, we introduce the vector notation $\boldsymbol{\zeta} \triangleq (\zeta_1, \zeta_2)^\top$ and recall the threshold vector $\boldsymbol{\mu}$ defined in Section~\ref{subsec:problem_formulation}.
Theorem~\ref{thm:general_optimality} then transforms the threshold optimization problem into a root-finding problem: the optimal thresholds are characterized by the system of equations $\boldsymbol{\zeta}(\boldsymbol{\mu}) = \mathbf{1}$.
This reformulation is useful for two reasons.
First, the indices $\zeta_i$ are computed directly from observable steady-state quantities $(z_0^*, z_{1i}^*, z_{2i}^*, q_i^*)$, obviating the need to compute implicit gradients $\nabla_{\boldsymbol{\mu}} \mathbf{x}^*$.
Second, the condition $\zeta_i = 1$ has a natural interpretation: at optimality, the marginal benefit of tightening the threshold for class $i$ (reducing abandonment, improving pickup speed) must exactly equal its marginal cost (consuming driver capacity, degrading service for class $j$).
Since the sensitivity analysis reveals $\partial \zeta_i / \partial \mu_{1i} < 0$, when $\zeta_i > 1$, the marginal benefit exceeds the marginal cost, indicating that the threshold $\mu_{1i}^*$ should be \emph{raised} to restore balance; when $\zeta_i < 1$, the threshold should be \emph{lowered}.

% =============================================================================
% SUBSECTION: Sensitivity Analysis
% Purpose: Characterize the Jacobian structure to motivate coordinated control
% =============================================================================
\subsection{Sensitivity Analysis and Numerical Properties}
\label{sec:sensitivity_analysis}

% =============================================================================
% PARAGRAPH: Motivation for numerical analysis
% Purpose: Explain why closed-form analysis is intractable
% =============================================================================
Theorem~\ref{thm:general_optimality} characterizes the optimal thresholds as roots of the system $\boldsymbol{\zeta}(\boldsymbol{\mu}) = \mathbf{1}$.
However, designing an algorithm to find these roots requires understanding the sensitivity of $\boldsymbol{\zeta}$ to changes in $\boldsymbol{\mu}$, namely the Jacobian matrix $\mathbf{J} = \nabla_{\boldsymbol{\mu}} \boldsymbol{\zeta}$.
Deriving closed-form expressions for $\mathbf{J}$ is analytically intractable because each steady-state component $(z_0^*, z_{1i}^*, z_{2i}^*, q_i^*)$ depends implicitly on the threshold vector $\boldsymbol{\mu}$ through the equilibrium equations.
To circumvent this difficulty, we characterize the Jacobian structure through extensive numerical experiments, which reveal consistent patterns.

% =============================================================================
% PARAGRAPH: Key findings - two structural properties
% Purpose: State the main empirical findings about the Jacobian
% =============================================================================
Figures~\ref{fig:zeta_W_sensitivity} and~\ref{fig:zeta_R_sensitivity} display the response surfaces of $\zeta_i^W$ and $\zeta_i^R$ to variations in the threshold pair $(\mu_{11}, \mu_{12})$.
These surfaces reveal two critical structural properties that persist across objectives and parameter regimes:

\begin{enumerate}
\item \textbf{Direct monotonicity:} 
$\partial \zeta_i / \partial \mu_{1i} < 0$.
Increasing the threshold $\mu_{1i}$ restricts matching opportunities for class $i$: fewer passengers are matched, the queue $q_i^*$ grows, and drivers accumulate in the idle pool $z_0^*$.
All three terms of $\zeta_i$ decrease as a result: the abandonment-reduction benefit shrinks (shorter pickup-to-queue ratio), the direct opportunity cost falls (more idle drivers available), and the cross-market externality weakens.
Thus, raising $\mu_{1i}$ mechanically reduces $\zeta_i$.

\item \textbf{Cross-market coupling:} 
$\partial \zeta_i / \partial \mu_{1j} > 0$ for $j \neq i$.
Raising the threshold for class $j$ (tightening matching) reduces the matching volume for class $j$, which alters the equilibrium in favor of class $i$.
As a result, $\zeta_i$ increases.
This positive cross-sensitivity means that threshold adjustments in one market spill over to affect the optimality conditions in the other, creating the coupling that necessitates coordinated control.
\end{enumerate}

\begin{figure}[!htbp]
    \centering
    \caption{Sensitivity analysis of Welfare-based Key Matching Indices ($\zeta^W$)}
    \label{fig:zeta_W_sensitivity}
    \begin{subfigure}[b]{0.48\textwidth}
        \centering
        \captionsetup{font=normalsize}
        \includegraphics[width=\textwidth]{figs/zeta_W_i1_3D.pdf}
        \caption{3D plot of $\zeta_1^W$}
    \end{subfigure}
    \hfill
    \begin{subfigure}[b]{0.42\textwidth}
        \centering
        \captionsetup{font=normalsize}
        \includegraphics[width=\textwidth]{figs/zeta_W_i1_contour.pdf}
        \caption{Contour of $\zeta_1^W$}
    \end{subfigure}
    
    \vspace{0.5em}
    \begin{subfigure}[b]{0.48\textwidth}
        \centering
        \captionsetup{font=normalsize}
        \includegraphics[width=\textwidth]{figs/zeta_W_i2_3D.pdf}
        \caption{3D plot of $\zeta_2^W$}
    \end{subfigure}
    \hfill
    \begin{subfigure}[b]{0.42\textwidth}
        \centering
        \captionsetup{font=normalsize}
        \includegraphics[width=\textwidth]{figs/zeta_W_i2_contour.pdf}
        \caption{Contour of $\zeta_2^W$}
    \end{subfigure}
    \begin{minipage}{0.9\textwidth}
    \footnotesize\raggedright
    \textit{Notes.} The plots show the response of $\zeta_1^W$ and $\zeta_2^W$ to variations 
    in matching thresholds $\mu_{11}$ and $\mu_{12}$. The 3D surfaces and contour plots 
    illustrate the direct and cross-market effects, confirming that $\zeta_i^W$ decreases 
    monotonically with its own threshold $\mu_{1i}$ and increases with the competing 
    threshold $\mu_{1j}$ ($j \neq i$).
    Base parameters: $\lambda_1=2.0$, $\lambda_2=3.0$, $\theta_{01}=2.0$, $\theta_{02}=1.0$, 
    $\theta_{11}=3.0$, $\theta_{12}=4.0$, $C=50.0$, $\alpha_1=\alpha_2=0.5$, $\mu_2=10.0$.
    \end{minipage}
\end{figure}

\begin{figure}[!htbp]
    \centering
    \caption{Sensitivity analysis of Revenue-based Key Matching Indices ($\zeta^R$)}
    \label{fig:zeta_R_sensitivity}
    \begin{subfigure}[b]{0.48\textwidth}
        \centering
        \captionsetup{font=normalsize}
        \includegraphics[width=\textwidth]{figs/zeta_R_i1_3D.pdf}
        \caption{3D plot of $\zeta_1^R$}
    \end{subfigure}
    \hfill
    \begin{subfigure}[b]{0.42\textwidth}
        \centering
        \captionsetup{font=normalsize}
        \includegraphics[width=\textwidth]{figs/zeta_R_i1_contour.pdf}
        \caption{Contour of $\zeta_1^R$}
    \end{subfigure}
    
    \vspace{0.5em}
    \begin{subfigure}[b]{0.48\textwidth}
        \centering
        \captionsetup{font=normalsize}
        \includegraphics[width=\textwidth]{figs/zeta_R_i2_3D.pdf}
        \caption{3D plot of $\zeta_2^R$}
    \end{subfigure}
    \hfill
    \begin{subfigure}[b]{0.42\textwidth}
        \centering
        \captionsetup{font=normalsize}
        \includegraphics[width=\textwidth]{figs/zeta_R_i2_contour.pdf}
        \caption{Contour of $\zeta_2^R$}
    \end{subfigure}
    \begin{minipage}{0.9\textwidth}
    \footnotesize\raggedright
    \textit{Notes.} Similar to the welfare case, $\zeta_i^R$ exhibits monotonic decay with respect to $\mu_{1i}$ and coupled increase with respect to $\mu_{1j}$, confirming the robustness of the Jacobian structure across different objectives.
    \end{minipage}
\end{figure}

% =============================================================================
% PARAGRAPH: Diagonal dominance and its implications
% Purpose: Establish the key property that enables algorithm design
% =============================================================================
\vspace{0.3em}

A crucial finding is that the \emph{direct effect dominates the cross-market effect}:
\(
\left| \frac{\partial \zeta_i}{\partial \mu_{1i}} \right|
>
\left| \frac{\partial \zeta_i}{\partial \mu_{1j}} \right|\quad\text{for all } i \neq j.
\)
This inequality ensures that the Jacobian matrix $\mathbf{J}$ is diagonally dominant, which has two important consequences.
First, diagonal dominance guarantees that $\mathbf{J}$ is invertible, so the system $\boldsymbol{\zeta}(\boldsymbol{\mu}) = \mathbf{1}$ admits a locally unique solution.
Second, it ensures the local stability of Newton-type iterations, preventing oscillatory or divergent behavior.

% =============================================================================
% PARAGRAPH: Why coordination is necessary
% Purpose: Explain the failure mode of naive control strategies
% =============================================================================
However, the non-zero off-diagonal entries underscore a critical limitation of independent tuning strategies.
Consider a naive coordinate-descent approach that adjusts $\mu_{11}$ and $\mu_{12}$ in isolation, ignoring cross-market effects.
Suppose $\zeta_1 > 1$ initially; the controller raises $\mu_{11}$ to reduce $\zeta_1$ toward unity.
Due to cross-market coupling ($\partial \zeta_2 / \partial \mu_{11} > 0$), this adjustment simultaneously raises $\zeta_2$.
If $\zeta_2$ now exceeds unity, the controller responds by raising $\mu_{12}$, which in turn raises $\zeta_1$ back above unity via $\partial \zeta_1 / \partial \mu_{12} > 0$.
This ping-pong dynamic can produce coupled oscillations that fail to converge.
The \emph{coordinated} control algorithm developed in the next subsection explicitly accounts for the full Jacobian structure to avoid such instabilities.

% =============================================================================
% SUBSECTION: Coordinated Control Algorithm
% Purpose: Present the Newton-Raphson algorithm for threshold optimization
% =============================================================================
\subsection{Coordinated Control Algorithm}
\label{subsec:coordinated_control}

% =============================================================================
% PARAGRAPH: Problem reformulation
% Purpose: Cast threshold optimization as root-finding
% =============================================================================
The stationarity characterization in Theorem~\ref{thm:general_optimality} reduces the threshold optimization problem to solving the nonlinear system
\begin{equation*}
  \boldsymbol{\zeta}(\boldsymbol{\mu}) = \mathbf{1}.
\end{equation*}
Standard root-finding algorithms can be applied, but the diagonal-dominant Jacobian structure established in Section~\ref{sec:sensitivity_analysis} suggests that Newton-type methods are particularly well-suited.

% =============================================================================
% PARAGRAPH: Newton-Raphson update
% Purpose: Present the core iteration formula
% =============================================================================
We employ a damped Newton--Raphson iteration that exploits the full Jacobian:
\begin{equation}
\boldsymbol{\mu}_{k+1}
=
\boldsymbol{\mu}_k
-
\eta\,
\hat{\mathbf{J}}^{-1}
\bigl(\boldsymbol{\zeta}(\boldsymbol{\mu}_k)-\mathbf{1}\bigr),
\label{eq:newton_update}
\end{equation}
where $\eta \in (0,1]$ is a step-size parameter that ensures stability, and $\hat{\mathbf{J}}$ denotes a numerical estimate of the Jacobian matrix $\mathbf{J} = \nabla_{\boldsymbol{\mu}} \boldsymbol{\zeta}$.
The damping parameter $\eta$ can be tuned via line search or set to a conservative value (e.g., $\eta = 0.5$) for robustness.

At each iteration, the algorithm estimates the steady-state variables by time-averaging, computes the current KMI values via~\eqref{eq:kmi_welfare} or~\eqref{eq:kmi_revenue}, estimates the Jacobian via finite differences, and applies the Newton update~\eqref{eq:newton_update} with projection onto $\boldsymbol{\mu} \geq \mathbf{0}$.
The complete procedure is given in Algorithm~\ref{alg:coordinated_control} in the Appendix.

\section{Numerical Experiments}\label{sec:numerical_experiments}

We conduct numerical experiments to validate the KMI-guided control, quantify the welfare-revenue frontier, and assess the value of coordination.
For concreteness, we refer to Class~1 as the \emph{premium} segment (e.g., Uber Comfort) and Class~2 as the \emph{regular} segment (e.g., UberX) throughout this section.
The analysis focuses on three canonical load regimes: (i) \textit{low-load}, where supply exceeds demand; (ii) \textit{balanced-load}, where queues coexist with idle capacity; and (iii) \textit{congested}, where demand outstrips supply.
Three main findings emerge: the KMI intersection reliably identifies the optimal operating point, the coordinated controller substantially outperforms decoupled and static benchmarks, and the welfare-revenue trade-off intensifies with congestion.
All quantitative results rely on steady-state statistics obtained via long-run tail averages.
Detailed simulation parameters and algorithmic tuning details are provided in Appendix~\ref{app:experiment_details}.

\subsection{Validation and Structural Geometry}
\label{sec:validate}

We first verify the optimality conditions derived in Section~\ref{sec:matching_policies}. 
First we define the First-Level Judgment (FLJ) index:
\begin{equation}
\mathrm{FLJ} \triangleq \sum_{i\in\{1,2\}} \frac{\alpha_{2} (z_{1i}^*/z_0^*)}{s_i}, \quad \text{where} \quad s_i = 1 - \alpha_{1}\frac{z_{1i}^*\theta_{1i}}{q_i^*\theta_{0i}}.
\label{eq:flj_def_exp}
\end{equation}
Theory posits that $\mathrm{FLJ}=1$ is necessary for optimality.
Figure~\ref{fig:flj-scatter} plots the welfare and revenue rates against the FLJ index.
While the performance envelope peaks near $\mathrm{FLJ}=1$, suboptimal configurations also satisfy this condition, confirming that the aggregate condition is necessary but not sufficient for optimality.

\begin{figure}[ht]
\centering
\caption{Necessity of the FLJ condition}
\label{fig:flj-scatter}
\begin{subfigure}{0.49\linewidth}
\centering
    \((a)\)\\
\includegraphics[width=\linewidth]{outputs/exp1/FLJsum_vs_welfare_rate_scatter_logX.pdf}
\end{subfigure}\hfill
\begin{subfigure}{0.49\linewidth}
\centering
    \((b)\)\\
\includegraphics[width=\linewidth]{outputs/exp1/FLJsum_vs_revenue_rate_scatter_logX.pdf}
\end{subfigure}
\\[0.4em]
\begin{subfigure}{0.78\linewidth}
\centering
    \((c)\)\\
\includegraphics[width=\linewidth]{outputs/exp1/FLJsum_vs_welfare_revenue_rate_scatter.pdf}
\end{subfigure}
\begin{subfigure}{0.9\textwidth}
    \footnotesize\raggedright
    \textit{Notes.} Vertical line marks $\mathrm{FLJ}=1$. High values concentrate 
    near this locus, yet suboptimal rates also occur at $\mathrm{FLJ}=1$, 
    highlighting necessity rather than sufficiency.
    \end{subfigure}
\end{figure}

\subsubsection{Objective Function Topography and Solution Structure}
To resolve the sufficiency gap and understand the solution geometry, we examine the topography of the objective in the $(\mu_{11},\mu_{12})$ plane using stochastic steady-state averages.
For welfare, Figure~\ref{fig:welfare-3d} shows a unimodal landscape with an isolated peak; the contour set is smooth and approximately convex.
Revenue displays a comparable geometry (Figure~\ref{fig:revenue-3d}), with a single global maximizer.

These plots also reveal pronounced anisotropy and curvature heterogeneity: slopes vary sharply across directions, and under high load, one observes ridges and plateaus.
Directly tuning $(\mu_{11},\mu_{12})$ is therefore difficult; the mapping from thresholds to the objective is nonlinear and regime-dependent.
The KMI structure addresses this difficulty.
By monitoring the state through $(\zeta_1,\zeta_2)$ and steering these indices toward one, control actions align with the unique maximizer and remain stable despite local irregularities.
Intuitively, while the objective surface in $(\mu_{11},\mu_{12})$ has complex curvature, the KMI provides a stable gradient field that guides the search toward the optimum.

\begin{figure}[ht]
\centering
\caption{Welfare objective topography (stochastic steady state)}
\label{fig:welfare-3d}
\begin{subfigure}{0.47\linewidth}
\centering
    \((a)\)\\
\includegraphics[width=\linewidth]{outputs/exp2_W/exp2_thresholds_busy_3D.pdf}
\end{subfigure}\hfill
\begin{subfigure}{0.47\linewidth}
\centering
    \((b)\)\\
\includegraphics[width=\linewidth]{outputs/exp2_W/exp2_thresholds_busy_contour.pdf}
\end{subfigure}
\begin{subfigure}{0.9\textwidth}
    \footnotesize\raggedright
    \textit{Notes.} (a) 3D surface of the threshold-based welfare proxy over $(\mu_{11},\mu_{12})$. 
    (b) Corresponding contour plot. The landscape has an isolated peak with smooth gradients yet pronounced anisotropy.
\end{subfigure}
\end{figure}

We next superimpose the KMI iso-structure.
Figures~\ref{fig:welfare-validate} and \ref{fig:revenue-validate} plot the $\zeta_1=1$ and $\zeta_2=1$ contours.
The numerical evidence shows that their unique intersection coincides with the objective peak in both cases.
This affirms the optimality condition and clarifies the operational role of KMIs as a state-dependent indicator: thresholds should be steered so that $(\zeta_1,\zeta_2)\to(1,1)$.
Given the surface complexity in $(\mu_{11},\mu_{12})$, blind tuning is fragile; KMI guidance regularizes search and delivers the maximizer reliably.
For completeness, the corresponding 3D landscapes of the KMI surfaces are provided in Appendix~\ref{app:3d_landscapes}.

\begin{figure}[ht]
\centering
\caption{Welfare validation: KMI intersection aligns with the optimum}
\label{fig:welfare-validate}
\centering
\includegraphics[width=0.6\linewidth]{outputs/exp2_W/exp2_kmi_busy_contour.pdf}
\begin{subfigure}{0.9\textwidth}
    \footnotesize\raggedright
    \textit{Notes.} Contours of $\zeta_1^W=1$ and $\zeta_2^W=1$ whose unique intersection aligns with the welfare-optimal thresholds.
\end{subfigure}
\end{figure}

\begin{figure}[ht]
\centering
\caption{Revenue objective topography (stochastic steady state)}
\label{fig:revenue-3d}
\begin{subfigure}{0.47\linewidth}
\centering
    \((a)\)\\
\includegraphics[width=\linewidth]{outputs/exp2_R/exp2_thresholds_busy_3D.pdf}
\end{subfigure}\hfill
\begin{subfigure}{0.47\linewidth}
\centering
    \((b)\)\\
\includegraphics[width=\linewidth]{outputs/exp2_R/exp2_thresholds_busy_contour.pdf}
\end{subfigure}
\begin{subfigure}{0.9\textwidth}
    \footnotesize\raggedright
    \textit{Notes.} (a) 3D surface of the threshold-based revenue proxy. (b) Corresponding contour plot.
\end{subfigure}
\end{figure}

\begin{figure}[ht]
\centering
\caption{Revenue validation: KMI intersection aligns with the optimum}
\label{fig:revenue-validate}
\centering
\includegraphics[width=0.6\linewidth]{outputs/exp2_R/exp2_kmi_busy_contour.pdf}
\begin{subfigure}{0.9\textwidth}
    \footnotesize\raggedright
    \textit{Notes.} Contours of $\zeta_1^R=1$ and $\zeta_2^R=1$ whose unique intersection aligns with the revenue-optimal thresholds.
\end{subfigure}
\end{figure}

\subsubsection{Quantitative Comparison of Optimal Policies}
We now quantify the trade-off between welfare and revenue maximization under the baseline parameterization.
Table~\ref{tab:exp-table} reports the optimal KMI-based policies and their corresponding steady-state performance metrics.

To maximize revenue, the platform shifts its policy from the welfare-optimal thresholds of $(\mu_{11}^*, \mu_{12}^*) \approx (9.99, 15.18)$ to a more permissive stance at $(10.58, 15.64)$.
This adjustment yields a marginal revenue gain of $0.14\%$ but comes at the cost of degraded service quality, as evidenced by increases in the premium and regular customer queues ($q_1^*$ and $q_2^*$) of $10.4\%$ and $4.1\%$.
The revenue-optimal policy achieves $99.8\%$ of the maximum possible welfare, indicating a tight welfare-revenue frontier in the balanced-load regime.

\begin{table}[ht]
\centering
\caption{Optimal thresholds and steady-state metrics (normalized by $K$; baseline regime)}
\label{tab:exp-table}
\small
\begin{tabular}{lccccccc}
\hline
Policy & $\mu_{11}^*$ & $\mu_{12}^*$ & Welfare rate & Revenue rate & $q_1^*$ & $q_2^*$ & $z_0^*$ \\
\hline
Welfare-opt & 9.99 & 15.18 & 3.512 & 4.894 & 0.087 & 0.230 & 0.490 \\
Revenue-opt & 10.58 & 15.64 & 3.506 & 4.901 & 0.096 & 0.239 & 0.497 \\
\hline
\end{tabular}
\vspace{0.25em}
\begin{minipage}{0.95\linewidth}
\footnotesize Notes: Rates are per unit time in fluid scale. Welfare rate is $\mu_2(z_{21}^*{+}z_{22}^*)$; 
revenue rate is $p_1\mu_2 z_{21}^*{+}p_2\mu_2 z_{22}^*$.
\end{minipage}
\end{table}

\subsection{Benchmarking and the Value of Coordination}
\label{sec:benchmarking}

We benchmark the KMI-based control against two rational alternatives: (i) a \textit{decoupled proxy} that optimizes thresholds sequentially, ignoring cross-segment externalities; and (ii) a \textit{static-priority} rule fixed at $(\mu_{11}, \mu_{12}) \approx (4,8)$.

Table~\ref{tab:benchmarks-congested} summarizes performance in the congested regime.
The KMI-based control dominates both benchmarks.
While the revenue gain over the decoupled proxy is marginal ($0.2\%$), the coordination gain in \textit{social welfare} is substantial ($2.9\%$).
This asymmetry arises because the decoupled proxy can approximately locate the revenue-maximizing point (since revenue is concentrated in the premium class), but it ignores the externality imposed on the regular class, leading to a larger welfare loss.
The primary value of the KMI framework therefore lies not merely in revenue, but in internalizing externalities to prevent inefficient cannibalization of the regular class.
By explicitly accounting for the shadow cost of capacity, the coordinated controller achieves a superior Pareto outcome.

We then quantify the \textbf{value of dynamism} by contrasting KMI control with the static-priority rule.
The state-dependent KMI policies deliver substantial gains: the welfare-optimal KMI increases welfare by $19.4\%$ and revenue by $8.7\%$, while the revenue-optimal KMI boosts revenue by $12.0\%$ and welfare by $17.2\%$.
These results underscore the importance of adapting matching criteria to real-time system state, especially under high load.

%The analysis also quantifies the welfare--revenue trade-off. To maximize profit, 
%the platform tightens the regular class threshold ($\mu_{12}$), reducing premium 
%queues ($q_1^*$) by $40.6\%$ at the cost of a $55.0\%$ increase in 
%regular queues ($q_2^*$) relative to the welfare-optimal policy.
The benchmarks also enrich our understanding of the welfare-revenue trade-off.
Under congestion, the revenue-optimal KMI policy reduces the premium queue ($q_1^*$) by $40.6\%$ relative to its welfare-optimal counterpart, but at the cost of a $55.0\%$ increase in the regular queue ($q_2^*$).
This trade-off shows how profit maximization reallocates scarce capacity toward the high-margin segment.
The KMI framework effectively manages this trade-off, substantially outperforming both the myopic decoupled policy and the non-adaptive static rule.

\begin{table}[htbp]
\centering
\caption{Performance benchmarking in the congested regime (normalized by $K$)}
\label{tab:benchmarks-congested}
\small
\begin{tabular}{lcccc}
\toprule
Metric & Welfare-Opt (KMI) & Revenue-Opt (KMI) & Decoupled & Static Priority \\
\midrule
Welfare Rate & \textbf{6.1003} & 5.9880 & 5.9275 & 5.1089 \\
Revenue Rate & 8.3013 & \textbf{8.5484} & 8.5328 & 7.6345 \\
Premium Queue ($q_1^*$) & 0.6216 & 0.3690 & 0.3219 & 0.0880 \\
Regular Queue ($q_2^*$) & 1.1334 & 1.7568 & 1.8142 & 2.7246 \\
Idle Fraction ($z_0^*$) & 0.2168 & 0.2188 & 0.1895 & 0.1456 \\
\bottomrule
\end{tabular}
\end{table}

\subsection{Dynamic Control and Convergence}
\label{sec:dynamic}

This section evaluates the performance of the $\zeta$-guided controller, focusing on convergence under time-varying demand.
Technical details regarding the convergence rate and hyperparameter tuning are deferred to Appendix~\ref{app:tuning}.

\subsubsection{Convergence under Time-Varying Demand}
We demonstrate convergence under non-stationary demand by initializing the system with suboptimal thresholds and subjecting it to piecewise-constant arrival rates that shift every 20 time units.
Figure~\ref{fig:dynamic} illustrates the response.
Matching thresholds $\mu_{1i}(t)$ adapt rapidly to load changes, converging to stable equilibria within each interval.
Concurrently, the KMIs $\zeta_i^W(t)$ are steered toward unity, stabilizing the welfare rate near its theoretical maximum.
This confirms that the controller effectively tracks the optimum under dynamic conditions.

\begin{figure}[htbp]
\centering
\caption{Dynamic performance under time-varying arrivals}
\label{fig:dynamic}
\begin{subfigure}{0.32\linewidth}
\centering
    \captionsetup{font=normalsize}
    \includegraphics[width=\linewidth]{outputs/exp3/exp3_mu_thresholds_timeseries.pdf}
    \subcaption{Thresholds $\mu_{1i}(t)$}
\end{subfigure}\hfill
\begin{subfigure}{0.32\linewidth}
\centering
    \captionsetup{font=normalsize}
    \includegraphics[width=\linewidth]{outputs/exp3/exp3_zetaW_timeseries.pdf}
    \subcaption{KMI $\zeta_i^W(t)$}
\end{subfigure}\hfill
\begin{subfigure}{0.32\linewidth}
\centering
    \captionsetup{font=normalsize}
    \includegraphics[width=\linewidth]{outputs/exp3/exp3_welfare_timeseries.pdf}
    \subcaption{Welfare Rate}
\end{subfigure}
\begin{subfigure}{0.95\textwidth}
    \footnotesize\raggedright
    \textit{Notes.} Arrival rates $(\lambda_1, \lambda_2)$ shift every 20 time units: 
    $(2.0, 1.0) \to (2.5, 2.0) \to (4.0, 4.0) \to (2.0, 2.0) \to (3.0, 4.0)$. 
    The controller maintains stability across transitions.
\end{subfigure}
\end{figure}

\subsubsection{Global Convergence}
%To assess robustness to initialization, we perform a multi-start experiment 
%from a grid of 81 distinct threshold pairs $(\mu_{11}, \mu_{12})$. 
%Figure~\ref{fig:exp3-multistart} depicts the trajectories. All paths 
%converge to a unique fixed point, which coincides with the revenue-optimal 
%policy identified in Section~\ref{sec:benchmarking}. This empirical evidence 
%supports the global convergence of the controller within the operational domain.

To assess the controller's robustness to initial conditions, we conduct a multi-start experiment.
The thresholds \((\mu_{11}, \mu_{12})\) are initialized from a grid of 81 distinct points, including those far from the optimum.
Figure~\ref{fig:exp3-multistart} shows the resulting trajectories in the threshold space and their corresponding time series.
All paths, regardless of their origin, converge to a single fixed point.
This empirically identified attractor coincides with the revenue-optimal policy reported in Section~\ref{sec:matching_policies}, providing strong evidence for the controller's global convergence property within the tested parameter space.

\begin{figure}[ht]
\centering
\caption{Multi-start convergence analysis}
\begin{subfigure}{0.6\linewidth}
\centering
\((a)\)\\
\includegraphics[width=\linewidth]{outputs/exp3/exp3_convergence_paths.pdf}
\end{subfigure}\hfill
\begin{subfigure}{0.8\linewidth}
\centering
\((b)\)\\
\includegraphics[width=\linewidth]{outputs/exp3/exp3_multistart_mu_timeseries.pdf}
\end{subfigure}
\begin{subfigure}{0.9\textwidth}
    \footnotesize\raggedright
    \textit{Notes.} (a) Phase Plane Trajectories. (b) Time Series.
	\end{subfigure}
\label{fig:exp3-multistart}
\end{figure}

\subsection{Sensitivity Analysis and Regime-Dependent Insights}
\label{sec:6-sensitivity}
We now quantify the sensitivity of the optimal policy to three key operational dimensions: aggregate system load, driver supply, and customer post-match patience.
The analysis relies on global policy scans over the discrete grid $(\mu_{11},\mu_{12})\in[1,20]^2$ within our stochastic simulation framework to identify the welfare- and revenue-maximizing policies under each parametric setting.

\subsubsection{Impact of Aggregate Load}
Table~\ref{tab:regime-sweep} documents the optimal policies and corresponding system performance across low, balanced, and high-load regimes.
A clear structural pattern emerges: the value of service differentiation is contingent on system congestion.
In the low-load regime, where driver supply is abundant relative to demand, queues are negligible, and the welfare- and revenue-optimal policies coincide.
As the system transitions to a high-load, congested state, the objectives diverge.
The revenue-maximizing policy intensifies service differentiation by substantially tightening the matching threshold for the regular class (a higher $\mu_{12}^*$).
This strategy prioritizes the high-margin premium segment, reducing its queue length ($q_1^*$) at the expense of a significantly longer queue for regular customers ($q_2^*$).
In contrast, the welfare-optimal policy adopts a more balanced approach, mitigating the negative externality imposed on the regular class.
The KMI-based framework provides the tool to precisely quantify the optimal degree of this trade-off, determining exactly how much service to the regular class must be curtailed to maximize revenue.
This result formally demonstrates how profit-seeking behavior reallocates scarce capacity in congested systems.

\input{outputs/generated/exp5_regime_table}

\subsubsection{Impact of Driver Supply}
Table~\ref{tab:K-sens} isolates the effect of driver supply, $K$, holding total customer arrival rates constant.
The results align with theoretical expectations for queueing systems.
As driver supply becomes scarcer (lower $K$), the system becomes more congested, leading to longer queues for both customer classes.
In response, the optimal policies become stricter, particularly for the lower-priority regular class under the revenue objective, to ration the limited service capacity more effectively.
Conversely, as driver supply increases, the system's congestion level eases, customer queues diminish, and the matching thresholds become more permissive.
The divergence between the welfare- and revenue-optimal policies narrows as supply grows, underscoring that the conflict between social good and platform profit is most acute when capacity is constrained.

\input{outputs/generated/exp5_K_table}

\subsubsection{Impact of Post-Match Patience}
We examine the role of post-match patience, parameterized by the abandonment hazard rates $(\theta_{11}, \theta_{12})$, where higher values signify less patient customers.
As shown in Table~\ref{tab:theta1-sens}, when customers are less patient, the optimal policy responds by tightening matching thresholds.
This strategic adjustment reduces average pickup times, thereby mitigating the risk of costly post-match abandonments.
This defensive posture, however, comes at the cost of lower overall welfare and revenue.
When customers are more patient (lower $\theta$ values), the system can tolerate longer pickup times, allowing for more relaxed thresholds and achieving superior performance.
Across all patience levels, the revenue-optimal policy is consistently more aggressive in restricting access for the regular class than its welfare-oriented counterpart, a pattern that is amplified when patience is low.

\input{outputs/generated/exp5_theta_table}

\subsection{Managerial Insights}
\label{sec:insights}
Our numerical investigation yields several insights for platform operations.
The KMI framework, validated as a reliable state-dependent indicator, provides the foundation for the control approach.

\begin{itemize}
    \item \textbf{Differentiation Driven by Scarcity:} 
    The optimal policy adapts to market conditions.
    Under scarcity (high load or low supply), the conflict between welfare and revenue widens, necessitating increased service differentiation to protect high-margin segments.
    Under abundance, objectives align, and differentiation diminishes.
    This formalizes the intuition that prioritization intensity should dynamically track resource scarcity.
    
    \item \textbf{Model-Based Control via KMI:}
    Targeting the $(\zeta_1,\zeta_2)=(1,1)$ condition provides a systematic threshold adjustment rule.
    This approach replaces reactive experimentation (such as A/B testing) with a model-based search that exploits the structure of the optimality conditions.
    Because the stationarity condition $\boldsymbol{\zeta} = \mathbf{1}$ defines a state-dependent efficient frontier rather than a fixed operating point, the optimal policy continuously adapts to changing load conditions, whereas static priority rules correspond to corner solutions that may be optimal under one demand pattern but severely suboptimal when conditions shift.
    
    \item \textbf{Stability and Externality Internalization:} 
    The gains from our control framework stem from two sources: (1) \emph{coordination}, by internalizing the cross-segment externalities (e.g., the impact of regular class admission on premium queues) that decoupled policies ignore, and (2) \emph{dynamism}, by adapting to the state-dependent system response.
\end{itemize}

\section{Conclusion}\label{sec:Conclusion}
% =============================================================================
% PARAGRAPH 1: SUMMARY OF CONTRIBUTIONS
% =============================================================================
% Role: Summarizes the paper's main problem, approach, and theoretical contributions.
% This is the "what we did" paragraph that recaps the key achievements.
%
% Internal Logic:
%   (1) Problem restatement: Reminds readers of the core problem addressed
%   (2) Modeling approach: Summarizes the methodological framework
%   (3) Central contribution: Highlights the main theoretical tool (KMI)
%   (4) Three-term structure: Emphasizes the decomposition and coupling (NEW)
%   (5) Key result: States the main theoretical finding (unity condition)
%   (6) Practical value: Explains the operational implications of the theory
% =============================================================================

% [Problem restatement]: Reminds readers of the core problem addressed
This paper addresses the control problem of managing heterogeneous demand classes within a shared-capacity ride-hailing network, where matching decisions for one class affect service quality for the other through endogenous pickup times that depend jointly on driver and passenger densities.
% [Modeling approach]: Summarizes the methodological framework used
By formulating a spatial queueing model with two-stage customer loss (pre-match abandonment and post-match cancellation), we established a tractable fluid framework that captures the endogenous feedback between matching decisions, idle driver density, and service quality.
% [Central contribution]: Highlights the main theoretical contribution (KMI)
The central theoretical contribution is the derivation of the Key Matching Indices (KMI), $(\zeta_1, \zeta_2)$, which transform complex state-dependent interactions into observable dimensionless metrics.
% [Three-term structure]: Emphasizes the decomposition and cross-market coupling (KEY NOVELTY)
Each index decomposes into three interpretable components: abandonment reduction, direct opportunity cost, and a cross-market term, with the third component explicitly coupling the two classes through the shared driver pool.
% [Key result]: States the main theoretical result about optimality conditions
We proved that driving both indices to unity characterizes the stationary points of the steady-state optimization problem for both welfare and revenue objectives, reducing a high-dimensional policy design problem to a tractable system of two coupled nonlinear equations.
% [Practical value]: Explains the operational implications of the theoretical result
This structural property provides an operational bridge between abstract fluid limits and real-time algorithmic control: the Jacobian-based update logic ensures convergence to optimal thresholds even under changing demand.

% =============================================================================
% PARAGRAPH 2: KEY FINDINGS AND INSIGHTS
% =============================================================================
% Role: Presents the main empirical findings and economic insights from the analysis.
% This is the "what we found" paragraph that highlights actionable insights.
%
% Internal Logic:
%   (1) Main finding: States the load-dependent nature of objective alignment
%   (2) Load regime clarification: Defines what we mean by load regimes
%   (3) Low-load result: Notes when objectives are aligned
%   (4) Congestion result: Describes what happens under resource constraints
%   (5) Mechanism explanation: Explains the capacity spillover effect
%   (6) Comparative performance: Contrasts with alternative approaches
%   (7) Validation: Notes the robustness of the optimization landscape
% =============================================================================

% [Main finding]: States the load-dependent nature of the welfare-revenue relationship
Our numerical investigations reveal that the alignment between social welfare and platform revenue is load-dependent.
% [Low-load result]: Notes when objectives are aligned
Under low load, both objectives can be optimized simultaneously with generous matching thresholds for both classes.
% [Congestion result]: Describes the trade-offs under resource constraints
Under congestion, however, revenue maximization requires disproportionate rationing of the regular class to preserve spatial density for premium users.
% [Mechanism explanation]: Explains the underlying capacity spillover mechanism
The divergence of optimal thresholds across objectives quantifies the high shadow cost of shared capacity when the two classes compete for idle drivers.
% [Comparative performance]: Contrasts KMI-guided control with alternatives
Compared to decoupled heuristics or static priority rules, the KMI-guided coordination explicitly internalizes these cross-market costs and delivers better outcomes across all load regimes tested.
% [Validation]: Notes the robustness demonstrated by the controller
The controller's ability to recover the optimum from any initial state further supports the well-behaved structure of the underlying objective.

% =============================================================================
% PARAGRAPH 3: LIMITATIONS AND FUTURE DIRECTIONS
% =============================================================================
% Role: Acknowledges the model's limitations and charts future research directions.
% This is the "what comes next" paragraph that positions future work.
%
% Internal Logic:
%   (1) Limitation acknowledgment: Notes the model's structural idealizations
%   (2) Future direction 1: Endogenous pricing and transient analysis
%   (3) Future direction 2: Multi-region network extensions
%   (4) Future direction 3: Learning-theoretic algorithmic enhancements
%   (5) Future direction 4: Fairness and equity constraints
%   (6) Closing statement: Summarizes the paper's broader contribution
% =============================================================================

% [Limitation acknowledgment]: Acknowledges the model's structural idealizations
Our model relies on several structural idealizations that suggest natural directions for future work.
% [Future direction 1]: Suggests relaxing pricing and arrival assumptions
First, our analysis presumes steady-state operation with exogenous pricing and arrivals; incorporating endogenous price sensitivity and short-horizon transients would bridge the gap between operational control and demand management.
% [Future direction 2]: Proposes extension to multi-region networks
Second, extending the framework to networked operations with multi-region imbalances would require characterizing the contraction properties of the KMI map in higher dimensions.
% [Future direction 3]: Suggests learning-theoretic algorithmic improvements
Third, deploying threshold controls in practice would benefit from learning-theoretic enhancements, such as confidence-adaptive step sizes, to handle model misspecification and data sparsity.
% [Future direction 4]: Notes the importance of incorporating fairness constraints
Finally, as platforms face increasing regulatory scrutiny regarding algorithmic fairness, integrating equity constraints directly into the KMI formulation would allow mapping a three-way frontier among welfare, revenue, and distributional fairness.

% [Closing statement]: Summarizes the paper's broader contribution to the field
Overall, this work provides a methodological foundation for moving from reactive, heuristic dispatching to model-based, state-dependent control in ride-hailing operations.

%\cite{Erlang1948,Dantzig1955,Dynkin1956,Bellman1957DP,Little1961,Skorokhod1961,McKean1965,Iglehart1965}
